\chapter{三维空间刚体运动}
\label{ch:rigid_body}

% ════════════════════════════════════════════════════════════════
%  第1章（数学基础部）：全量吸收 十四讲 ch3 + Barfoot ch6（内部编号7.x）+ SLAM Handbook ch2/Notation
%  自包含、主要证明内联、右扰动约定的前置。范围=旋转/位姿的【表示与互转 + 运动学】，
%  群/李代数/指数映射/扰动求导/伴随留给 \cref{ch:lie}。
% ════════════════════════════════════════════════════════════════

\begin{practice}[前置自测]
开始之前，先试答下列问题。能流畅作答可速读；卡住之处正是本章要解决的。每题后给“答错→回看”指引。
\begin{enumerate}
  \item “向量 $\mathbf{a}$” 与 “$\mathbf{a}$ 的坐标 $[a_1,a_2,a_3]^\top$” 是同一回事吗？后者依赖什么？（→ \cref{sec:rb-vector}）
  \item 旋转矩阵有 $9$ 个元素，但旋转只有 $3$ 个自由度。多出来的约束是什么？如何证明旋转矩阵正交？（→ \cref{sec:rb-rotation}）
  \item 已知“相机到世界”的位姿 $\mathbf{T}_{wc}$，平移部分 $\mathbf{t}_{wc}$ 的几何含义是什么？它等于 $-\mathbf{t}_{cw}$ 吗？（→ \cref{sec:rb-convention}）
  \item 为什么不能用 $3$ 个欧拉角无奇异地表示所有三维旋转？“万向锁”锁住了什么？（→ \cref{sec:rb-euler}）
  \item 四元数 $\mathbf{q}$ 和 $-\mathbf{q}$ 表示同一个旋转吗？Hamilton 与 JPL 约定的差别是什么？（→ \cref{sec:rb-quat}）
\end{enumerate}
\end{practice}

\section*{本章目标}
学完本章，你应当能够：
\begin{enumerate}
  \item \textbf{区分}向量与其坐标，熟练运用内积、外积与反对称算子 $(\cdot)^\wedge$ 及其\emph{全套恒等式}；
  \item \textbf{推导}旋转矩阵并\emph{证明} $\SO(3)$ 的正交约束、方向余弦含义、主动/被动两种解读；
  \item \textbf{运用}齐次坐标与变换矩阵 $\SE(3)$ 完成位姿复合、求逆与链式相消；
  \item \textbf{厘清}本书 / Barfoot / Handbook 三套坐标系记号，并避开 $\mathbf{t}_{12}\neq-\mathbf{t}_{21}$ 之类陷阱；
  \item \textbf{证明}罗德里格斯公式（几何分解法 + 矩阵指数级数法两向），并在旋转矩阵、轴角、欧拉角、四元数、Euler 参数、Gibbs 向量之间互转；
  \item \textbf{解释}欧拉角的万向锁与“三参数必有奇异”（Stuelpnagel 1964）的一般性结论；
  \item \textbf{写出}旋转/位姿运动学（Poisson 方程、哥氏加速度、Frenet--Serret、独轮车模型、欧拉角速率映射），为 \cref{ch:lie} 的李代数埋下伏笔；
  \item \textbf{建立}（选读）单刚体在 $\SE(3)$ 上的动力学——D'Eleuterio 广义质量阵 $\boldsymbol{\mathcal{M}}$、统一的牛顿--欧拉方程及其线性化，作为含动力学运动模型的过程方程；
  \item \textbf{使用} Eigen 几何模块完成旋转/位姿的表示、互转与坐标变换，并复现一道完整数值例。
\end{enumerate}

\subsection*{本章知识导航}
本章只回答一件事：\emph{如何用数学语言\textbf{表示}三维刚体的位置与朝向，以及各表示如何互转、如何随时间演化}。“如何\textbf{估计}它”留给 \cref{ch:lie}。结构如下：

\begin{center}
\small
\begin{tikzpicture}[
  node distance=5mm and 7mm, >=Stealth,
  blk/.style={draw, rounded corners, align=center, font=\small, inner sep=3pt, fill=main!6, draw=main!60!black},
  rep/.style={draw, rounded corners, align=center, font=\small, inner sep=3pt, fill=second!6, draw=second!60!black},
  ln/.style={->, thick, gray!60!black}]
\node[blk] (frame) {点·向量·坐标系\\ $(\cdot)^\wedge$};
\node[blk, right=of frame] (rot) {旋转矩阵\\ $\SO(3)$};
\node[blk, right=of rot] (pose) {齐次坐标·位姿\\ $\SE(3)$};
\node[rep, below=9mm of frame] (axis) {轴角\\ Rodrigues};
\node[rep, right=of axis] (euler) {欧拉角\\ 万向锁};
\node[rep, right=of euler] (quat) {四元数\\ Hamilton};
\node[blk, below=9mm of euler, fill=third!8, draw=third!60!black] (kin) {运动学\\ Poisson};
\draw[ln] (frame)--(rot); \draw[ln] (rot)--(pose);
\draw[ln] (rot.south) -- ++(0,-5mm) -| (axis.north);
\draw[ln] (axis)--(euler); \draw[ln] (euler)--(quat);
\draw[ln] (rot) to[bend left=10] (quat.north);
\draw[ln] (euler.south) -- (kin.north);
\end{tikzpicture}
\end{center}

\noindent\textbf{推荐路径（脉络层）}：\emph{骨干}=第 1--4 节（向量→旋转矩阵→位姿→记号约定）和第 5、7 节（轴角/Rodrigues、四元数）——后续\emph{所有}章节都依赖它们；\emph{次优}=第 6 节（欧拉角，理解奇异即可）、第 8 节（对比选型）、运动学（\cref{sec:rb-kinematics}，承接 \cref{ch:lie}）；\emph{选读}=单刚体动力学（\cref{sec:rb-dynamics}）、射影变换（\cref{sec:rb-projective}，服务相机模型章）、第 6.4 节的 Gibbs 向量。末节（\cref{sec:rb-eigen}）是 Eigen 落地。带（选读）标记的小节为选读。难度轴与脉络层正交：第 5、7 节属骨干但难度 $\bigstar\bigstar\bigstar$。

\subsection*{全书脉络定位：本章承上启下}
在进入正文之前，先把本章放回全书的主线里，看清它“从哪里来、到哪里去”。前面的章节为我们准备好了通用的数学语言与记号（向量、矩阵、群的初步约定，见 \nameref{ch:notation}），并交代了视觉 SLAM 的整体框架——一个机器人要完成定位与建图，第一步永远是\emph{回答“我在哪、朝哪看”}。但“在哪、朝哪看”这句自然语言无法直接计算，本章的任务正是把它\textbf{翻译成无歧义的数学对象}：用旋转矩阵、轴角、欧拉角、四元数表示\emph{朝向}，用平移向量表示\emph{位置}，再用位姿矩阵把二者合二为一。

\textbf{本章解决什么、不解决什么}：本章只回答\emph{表示}与\emph{互转}（以及它们如何随时间演化的运动学），\emph{不}回答“如何从含噪声的观测里估计出这个位姿”。后者需要把旋转放到流形上做微分与优化，那是紧随其后的 \cref{ch:lie}（李群与李代数）的主题——而 \cref{ch:lie} 的每一块基石（$\SO(3)/\SE(3)$ 的非线性约束、Rodrigues 公式、Poisson 方程、广义速度）都在本章先行铺好。再往后，\cref{ch:camera}（相机模型）会把本章的齐次坐标与位姿用于重投影，\cref{ch:vo}（视觉里程计）会反过来\emph{求解}位姿，\cref{ch:vio}（VIO）则直接调用本章的四元数运动学与哥氏加速度做 IMU 积分。可以说，本章是全书几何部分的\emph{地基}：读懂它，后续才不会在“坐标系方向写反、平移直接取负、欧拉角在大角度处崩溃”这些最常见的坑里反复栽跟头。

\subsection*{前置知识桥接}
本章只需线性代数基础（向量、矩阵、正交、特征值、迹）。\textbf{动机}：要让机器人“知道自己在哪、朝哪看”，必须先有一套\emph{无歧义}的数学语言描述刚体的位置与朝向；本章建立这套语言，后续\emph{所有}章节（估计、优化、感知、SLAM）都在其上展开。Barfoot\cite{barfoot2024state} 把刚体的“位置 + 朝向”合称\textbf{位姿（pose）}，数学上 $6$ 个自由度（$3$ 平移 $+$ $3$ 旋转）。

\subsection*{如果跳过本章会怎样}
\begin{itemize}
  \item 读 \cref{ch:lie} 时，你会困惑“为什么旋转不能直接相加、要搬到李代数上”——而那正源于本章 $\SO(3)$ 的非线性约束；
  \item 读 \cref{ch:camera}（相机模型）时，重投影 $\mathbf{z}=\pi(\mathbf{T}_{cw}\mathbf{p}_w)$ 里的位姿、齐次坐标、帧下标若没在这里厘清，几乎必然把坐标系方向写反；
  \item 读 \cref{ch:vio}（VIO）时，IMU 的姿态积分用的就是本章的 Poisson 方程与四元数运动学，不懂这里就无法理解预积分。
\end{itemize}

\begin{table}[H]\centering\small
\caption{本章阅读方式建议}
\begin{tabular}{lll}
\toprule
阅读方式 & 时间 & 适合谁 \\
\midrule
精读（含全部推导与练习） & 5--7 小时 & 首次系统学习 \\
速读（跳过冗长推导细节） & 2 小时 & 有基础、想对齐本书记号约定 \\
速查（只看小结的符号表 + 速查表） & 10 分钟 & 写代码时回来查公式 \\
\bottomrule
\end{tabular}
\end{table}

\begin{note}
\textbf{本章记号与约定.}\;
旋转矩阵记 $\mathbf{R}\in\SO(3)$（Barfoot\cite{barfoot2024state} 用 $\mathbf{C}$）；位姿 $\mathbf{T}\in\SE(3)$。
\textbf{帧下标“从右读到左、左为目标”}：$\mathbf{R}_{wc}$ 把\emph{相机系}坐标变到\emph{世界系}，$\mathbf{p}_w=\mathbf{R}_{wc}\mathbf{p}_c$；位姿同理 $\mathbf{T}_{wc}$。
\textbf{四元数取 Hamilton 约定}（$\mathrm{ij}=\mathrm{k}$，与 Eigen 一致），数学书写\emph{实部在前} $\mathbf{q}=[w,\mathbf{v}]^\top$（注意 Eigen 内存按 $(x,y,z,w)$ 存）。
反对称算子 $(\cdot)^\wedge:\mathbb{R}^3\to\mathbb{R}^{3\times3}$，逆为 $(\cdot)^\vee$。本章只讲\textbf{表示、互转与运动学}；群结构、李代数、指数映射、扰动与求导、伴随见 \cref{ch:lie}。
\textbf{字母隔离声明}：本章正文中 $\bigstar$ 仅作难度标记（$\bigstar$ 必学 / $\bigstar\bigstar$ 核心 / $\bigstar\bigstar\bigstar$ 进阶），与数学记号无关；轴角转角统一记 $\theta$，旋转李代数向量在 \cref{ch:lie} 记 $\boldsymbol{\phi}$。
\end{note}

% ════════════════════════════════════════════════════════════════
\section{点、向量与坐标系}
\label{sec:rb-vector}
\noindent\emph{节首导读：骨干必跟。本节建立全章最底层的概念（向量 vs 坐标）与最常用的工具（反对称算子），后续每一节都用到。}

\paragraph{动机：把“朝前方”变成数。}
“相机在 $(0,0,0)$、朝正前方”这种自然语言无法计算。三维空间由 $3$ 个轴组成，一个空间点的位置由 $3$ 个坐标指定；但刚体不光有\emph{位置}，还有\emph{姿态}（朝向），相机正可看作三维空间中的一个刚体\cite{gaoxiang2019slam14}。要让计算机处理刚体运动，必须用\emph{坐标}。但坐标之前，先要分清两个常被混为一谈的概念。

\paragraph{向量 $\neq$ 坐标（核心要点）。}
\emph{点}是空间的基本元素，没有长度、没有体积；连接两点的箭头是\emph{向量}。向量是空间中客观存在的几何量（带长度和方向），它\textbf{不依赖}任何坐标系，不需要与一组实数相关联——这是全章最易被忽略、却最该先立稳的区分\cite{gaoxiang2019slam14}。只有在指定一组\emph{基}（线性空间的一组基底）$(\mathbf{e}_1,\mathbf{e}_2,\mathbf{e}_3)$ 后，才能谈它的\emph{坐标}：
\begin{equation}
  \mathbf{a}=[\mathbf{e}_1,\mathbf{e}_2,\mathbf{e}_3]\begin{bmatrix}a_1\\a_2\\a_3\end{bmatrix}=a_1\mathbf{e}_1+a_2\mathbf{e}_2+a_3\mathbf{e}_3.
  \label{eq:coord}
\end{equation}
坐标 $[a_1,a_2,a_3]^\top$ 同时取决于\emph{向量本身}与\emph{基的选取}。同一向量在不同坐标系下坐标不同——这正是后面“坐标变换”要处理的。

\paragraph{右手系与左手系。}坐标系通常由 $3$ 个正交坐标轴组成。给定 $x$、$y$ 轴后，$z$ 轴由 $\mathbf{x}\times\mathbf{y}$（右手法则）或 $\mathbf{y}\times\mathbf{x}$（左手法则）确定，分别得右手系与左手系，二者第三轴方向相反。常见右手系库：OpenGL、3D Max；左手系库：Unity、Direct3D\cite{gaoxiang2019slam14}。本书除非特别声明，一律用\textbf{右手系}。

\paragraph{Barfoot 的 vectrix 记号（更严谨的“向量/坐标”分离）。}十四讲用文字强调“向量 $\neq$ 坐标”，Barfoot\cite{barfoot2024state} 则把它\emph{代数化}：把一个参考系 $\underline{\mathcal{F}}_1$ 的三个单位正交右手基向量 $\underline{1}_1,\underline{1}_2,\underline{1}_3$ 堆成一个 $3\times1$ 的“基向量列”——称为 \textbf{vectrix} $\underline{\mathcal{F}}_1=[\underline{1}_1;\underline{1}_2;\underline{1}_3]$，于是物理向量 $\underline{r}$ 与其在该系下的坐标列 $\mathbf{r}_1$ 严格分离：
\begin{equation}
  \underline{r}=r_1\underline{1}_1+r_2\underline{1}_2+r_3\underline{1}_3
  =\underbrace{[\underline{1}_1\ \underline{1}_2\ \underline{1}_3]}_{\underline{\mathcal{F}}_1^\top}\,\mathbf{r}_1
  =\mathbf{r}_1^\top\underline{\mathcal{F}}_1 .
\end{equation}
\textbf{下标数字 = 在哪个系里表达}。这套记号最大的好处：点积、叉积等运算可以“连基带坐标”一起算，化简时基的正交性会自动“吐出”单位阵——下面两小节就会看到。两书视角并列：十四讲更直觉（文字提醒），Barfoot 更严谨（记号上不可能混淆）。

\paragraph{内积。}两向量 $\mathbf{a},\mathbf{b}\in\mathbb{R}^3$，内积（点积）度量投影：
\begin{equation}
  \mathbf{a}\cdot\mathbf{b}=\mathbf{a}^\top\mathbf{b}=\sum_{i=1}^3 a_ib_i=|\mathbf{a}||\mathbf{b}|\cos\langle\mathbf{a},\mathbf{b}\rangle,
  \label{eq:rb-dot}
\end{equation}
其中 $\langle\mathbf{a},\mathbf{b}\rangle$ 是两向量夹角。用 vectrix 推一遍可看清“为什么坐标内积就是物理内积”：设 $\underline{r}=\mathbf{r}_1^\top\underline{\mathcal{F}}_1$、$\underline{s}=\underline{\mathcal{F}}_1^\top\mathbf{s}_1$，则
\begin{equation}
  \underline{r}\cdot\underline{s}
  =\mathbf{r}_1^\top\Big(\underline{\mathcal{F}}_1\cdot\underline{\mathcal{F}}_1^\top\Big)\mathbf{s}_1
  =\mathbf{r}_1^\top\,\mathbf{I}\,\mathbf{s}_1=\mathbf{r}_1^\top\mathbf{s}_1,
\end{equation}
中间矩阵 $\underline{\mathcal{F}}_1\cdot\underline{\mathcal{F}}_1^\top$ 的 $(i,j)$ 元是 $\underline{1}_i\cdot\underline{1}_j=\delta_{ij}$（基正交归一），故为单位阵。结论：内积只与两向量本身有关，\emph{与坐标系选取无关}——这正是“向量是客观对象”的代数体现。

\paragraph{外积与反对称算子 $(\cdot)^\wedge$。}外积结果是\emph{向量}，方向垂直于二者所张平面、大小 $|\mathbf{a}||\mathbf{b}|\sin\langle\mathbf{a},\mathbf{b}\rangle$ 等于两向量张成的平行四边形的有向面积。它可写成行列式展开，进而写成一个矩阵乘法，由此引入贯穿全书的\textbf{反对称算子} $(\cdot)^\wedge$：
\begin{equation}
  \mathbf{a}\times\mathbf{b}
  =\begin{vmatrix}\mathbf{e}_1&\mathbf{e}_2&\mathbf{e}_3\\a_1&a_2&a_3\\b_1&b_2&b_3\end{vmatrix}
  =\begin{bmatrix}a_2b_3-a_3b_2\\a_3b_1-a_1b_3\\a_1b_2-a_2b_1\end{bmatrix}
  =\underbrace{\begin{bmatrix}0&-a_3&a_2\\a_3&0&-a_1\\-a_2&a_1&0\end{bmatrix}}_{\mathbf{a}^\wedge}\mathbf{b}.
  \label{eq:cross}
\end{equation}
$(\cdot)^\wedge$ 把向量\emph{一一映射}为反对称（skew-symmetric）矩阵：任意向量唯一对应一个反对称阵，反之亦然；它的逆映射记 $(\cdot)^\vee$，即 $(\mathbf{a}^\wedge)^\vee=\mathbf{a}$。把外积写成 $\mathbf{a}^\wedge\mathbf{b}$ 的关键意义是：\textbf{外积变成了线性运算}——这使得后面无数推导（旋转点的求导、四元数到矩阵、Rodrigues、运动学）都能用矩阵代数机械地进行。用 vectrix 同样可导出 $(\cdot)^\wedge$：由 $\underline{1}_1\times\underline{1}_2=\underline{1}_3$ 等右手关系展开 $\underline{r}\times\underline{s}=\underline{\mathcal{F}}_1^\top\mathbf{r}_1^\wedge\mathbf{s}_1$，中间矩阵正是 \cref{eq:cross} 的 $\mathbf{a}^\wedge$\cite{barfoot2024state}。

\begin{note}
\textbf{记号源流.}\;十四讲\cite{gaoxiang2019slam14} 记 $\mathbf{a}^\wedge$；Barfoot 前半用 $\mathbf{a}^\times$、自其“机器人学约定”一节（源 §7.3.2）起改用 $\mathbf{a}^\wedge$；亦有文献写 $[\mathbf{a}]_\times$、$\hat{\mathbf{a}}$、$-[[\mathbf{a}]]$。Handbook\cite{carlone2026handbook} 用 $\boldsymbol{\theta}^\wedge$。本书统一 $(\cdot)^\wedge$。
\end{note}

\begin{insight}[反对称算子是“把几何写成线性代数”的钥匙]
外积本是一个几何运算（垂直、面积、右手），引入 $(\cdot)^\wedge$ 后它变成矩阵乘法。这不是雕虫小技：旋转的微分（\cref{eq:poisson}）、四元数到旋转矩阵（\cref{eq:quat2R}）、Rodrigues 的级数证明（\cref{thm:rodrigues-rb}）、扰动求导（\cref{ch:lie}）全都靠把 $\times$ 换成 $(\cdot)^\wedge$ 才得以机械推进。先把它的恒等式背熟，后面省一半力气。
\end{insight}

\paragraph{反对称矩阵恒等式全套（贯穿全书，务必熟记）。}下列恒等式十四讲与 Barfoot 散见于各处求迹/化简，这里集中证明并编号，后文反复引用。设 $\mathbf{a},\mathbf{b}\in\mathbb{R}^3$，$\mathbf{a}$ 在需要处为单位向量（$\mathbf{a}^\top\mathbf{a}=1$），$\mathbf{R}\in\SO(3)$。
\begin{align}
  \text{(I1)}\quad& \mathbf{a}^\wedge\mathbf{b}=-\mathbf{b}^\wedge\mathbf{a}, \label{eq:rb-id1}\\
  \text{(I2)}\quad& \mathbf{a}^\wedge\mathbf{a}=\mathbf{0},\qquad (\mathbf{a}^\wedge)^\top=-\mathbf{a}^\wedge,\qquad \operatorname{tr}(\mathbf{a}^\wedge)=0, \label{eq:rb-id2}\\
  \text{(I3)}\quad& \mathbf{a}^\wedge\mathbf{a}^\wedge=\mathbf{a}\mathbf{a}^\top-(\mathbf{a}^\top\mathbf{a})\mathbf{I}, \label{eq:rb-id3}\\
  \text{(I4)}\quad& \mathbf{a}^\wedge\mathbf{a}^\wedge\mathbf{a}^\wedge=-(\mathbf{a}^\top\mathbf{a})\,\mathbf{a}^\wedge, \label{eq:rb-id4}\\
  \text{(I5)}\quad& (\mathbf{R}\mathbf{a})^\wedge=\mathbf{R}\,\mathbf{a}^\wedge\mathbf{R}^\top, \label{eq:rb-id5}\\
  \text{(I6)}\quad& \mathbf{a}\times(\mathbf{b}\times\mathbf{c})=\mathbf{b}(\mathbf{a}^\top\mathbf{c})-\mathbf{c}(\mathbf{a}^\top\mathbf{b})\quad(\text{BAC--CAB}). \label{eq:rb-id6}
\end{align}
(I1) 即外积反交换；(I2) 把向量代入定义直接验证（对角线为零，转置变号，迹为零）。(I3) 是后文求迹的关键：逐元素展开 $\mathbf{a}^\wedge\mathbf{a}^\wedge$ 与 $\mathbf{a}\mathbf{a}^\top-(\mathbf{a}^\top\mathbf{a})\mathbf{I}$ 相等即可；由它立得当 $\mathbf{a}$ 为单位向量时 $\operatorname{tr}(\mathbf{a}^\wedge\mathbf{a}^\wedge)=\operatorname{tr}(\mathbf{a}\mathbf{a}^\top)-3\,\mathbf{a}^\top\mathbf{a}=1-3=-2$。(I4) 由 (I3) 右乘 $\mathbf{a}^\wedge$ 并用 (I2) 的 $\mathbf{a}^\top\mathbf{a}^\wedge=\mathbf{0}^\top$（即 $(\mathbf{a}^\wedge\mathbf{a})^\top=\mathbf{0}^\top$）得到。(I5) 与 (I6) 是 Barfoot 习题 7.3 / 三重积公式，证明见练习与 \cref{sec:rb-appendix}。

\begin{pitfall}[证明陷阱：把“正交”当成“旋转”，漏掉 $\det=+1$]
$\mathbf{R}^\top\mathbf{R}=\mathbf{I}$ 只保证 $\mathbf{R}\in\mathrm{O}(3)$（正交群），它含 $\det=+1$（真旋转 $\SO(3)$）与 $\det=-1$（含镜像/反射）两支。漏掉 $\det=+1$ 会把镜像当旋转：在 $\SO(3)$ 上做指数映射/插值时混入反射会导致不连续、不可积。\textbf{凡声称“旋转”，必须同时写 $\mathbf{R}^\top\mathbf{R}=\mathbf{I}$ 且 $\det\mathbf{R}=+1$。}
\end{pitfall}

\begin{practice}
\begin{enumerate}
  \item （证明）由定义验证 (I1)(I2)；并由 (I3) 推出 (I4)。\emph{答案线索}：(I4) 用 $\mathbf{a}^\wedge\mathbf{a}=\mathbf{0}$。
  \item （证明）用 $(\cdot)^\wedge$ 把三重积写成矩阵形式，证明 (I6) $\mathbf{a}\times(\mathbf{b}\times\mathbf{c})=\mathbf{b}(\mathbf{a}^\top\mathbf{c})-\mathbf{c}(\mathbf{a}^\top\mathbf{b})$。
  \item （推导）逐元素验证 (I3)：写出 $\mathbf{a}^\wedge\mathbf{a}^\wedge$ 的 $(1,1)$ 元 $-a_2^2-a_3^2$ 与 $\mathbf{a}\mathbf{a}^\top-(\mathbf{a}^\top\mathbf{a})\mathbf{I}$ 的 $(1,1)$ 元 $a_1^2-(a_1^2+a_2^2+a_3^2)$ 相等。
\end{enumerate}
\end{practice}

% ════════════════════════════════════════════════════════════════
\section{旋转矩阵}
\label{sec:rb-rotation}
\noindent\emph{节首导读：骨干必跟。旋转矩阵是“作用于点、做坐标变换”的主力表示；其正交性证明是本章第一处承重推导。}

\paragraph{动机：刚体运动 = 旋转 + 平移。}
把手机抛向空中，落地前它只改变\emph{位置}与\emph{朝向}，自身长度、各面夹角不变——这种保长度、保夹角的运动叫\textbf{欧氏变换}（Euclidean transform）\cite{gaoxiang2019slam14}。两个坐标系之间的运动恰好是一个旋转加一个平移。平移好办（坐标相加即可），难点在旋转。

\paragraph{由基变换导出旋转矩阵（几何路线）。}
设单位正交基 $(\mathbf{e}_1,\mathbf{e}_2,\mathbf{e}_3)$ 旋转成 $(\mathbf{e}_1',\mathbf{e}_2',\mathbf{e}_3')$，同一向量 $\mathbf{a}$（未随坐标系旋转而运动）在两组基下坐标为 $[a_1,a_2,a_3]^\top$ 与 $[a_1',a_2',a_3']^\top$。因向量本身没变，由坐标定义：
\begin{equation}
  [\mathbf{e}_1,\mathbf{e}_2,\mathbf{e}_3]\begin{bmatrix}a_1\\a_2\\a_3\end{bmatrix}
  =[\mathbf{e}_1',\mathbf{e}_2',\mathbf{e}_3']\begin{bmatrix}a_1'\\a_2'\\a_3'\end{bmatrix}.
  \label{eq:rb-samevec}
\end{equation}
两边\emph{左乘} $[\mathbf{e}_1,\mathbf{e}_2,\mathbf{e}_3]^\top=[\mathbf{e}_1^\top;\mathbf{e}_2^\top;\mathbf{e}_3^\top]$。因 $(\mathbf{e}_1,\mathbf{e}_2,\mathbf{e}_3)$ 单位正交，$\mathbf{e}_i^\top\mathbf{e}_j=\delta_{ij}$，左边系数矩阵化为单位阵，于是
\begin{equation}
  \begin{bmatrix}a_1\\a_2\\a_3\end{bmatrix}
  =\underbrace{\begin{bmatrix}\mathbf{e}_1^\top\mathbf{e}_1'&\mathbf{e}_1^\top\mathbf{e}_2'&\mathbf{e}_1^\top\mathbf{e}_3'\\
  \mathbf{e}_2^\top\mathbf{e}_1'&\mathbf{e}_2^\top\mathbf{e}_2'&\mathbf{e}_2^\top\mathbf{e}_3'\\
  \mathbf{e}_3^\top\mathbf{e}_1'&\mathbf{e}_3^\top\mathbf{e}_2'&\mathbf{e}_3^\top\mathbf{e}_3'\end{bmatrix}}_{\mathbf{R}}
  \begin{bmatrix}a_1'\\a_2'\\a_3'\end{bmatrix}.
  \label{eq:R-derive}
\end{equation}
矩阵 $\mathbf{R}$ 的每个元素是两组基向量的内积。因基向量长度为 $1$，内积即两基向量\emph{夹角的余弦}，故 $\mathbf{R}$ 又称\textbf{方向余弦矩阵}（direction cosine matrix, DCM）。它\emph{只取决于旋转本身}（只要旋转一样，$\mathbf{R}$ 就一样），刻画了同一向量在旋转前后坐标的变换关系。

\paragraph{Barfoot 的 vectrix 路线（同一结论的严谨版）。}两个共原点的系 $\underline{\mathcal{F}}_1,\underline{\mathcal{F}}_2$，同一向量 $\underline{r}=\underline{\mathcal{F}}_1^\top\mathbf{r}_1=\underline{\mathcal{F}}_2^\top\mathbf{r}_2$。两边点乘 $\underline{\mathcal{F}}_2$：
\begin{equation}
  \underline{\mathcal{F}}_2\cdot\underline{\mathcal{F}}_2^\top\,\mathbf{r}_2=\underline{\mathcal{F}}_2\cdot\underline{\mathcal{F}}_1^\top\,\mathbf{r}_1
  \;\Longrightarrow\;
  \mathbf{r}_2=\mathbf{C}_{21}\mathbf{r}_1,\qquad
  \mathbf{C}_{21}:=\underline{\mathcal{F}}_2\cdot\underline{\mathcal{F}}_1^\top,
\end{equation}
（用了 $\underline{\mathcal{F}}_2\cdot\underline{\mathcal{F}}_2^\top=\mathbf{I}$）。$\mathbf{C}_{21}$ 的 $(i,j)$ 元是 $\underline{2}_i\cdot\underline{1}_j$，同样是方向余弦。两书结论一致：$\mathbf{R}\leftrightarrow\mathbf{C}_{21}$，本书用 $\mathbf{R}$。

\begin{definition}[特殊正交群 $\SO(3)$]\label{def:so3-rep}
三维旋转矩阵的集合
\begin{equation}
  \SO(3)=\big\{\mathbf{R}\in\mathbb{R}^{3\times3}\,\big|\,\mathbf{R}\mathbf{R}^\top=\mathbf{I},\ \det\mathbf{R}=1\big\}.
  \label{eq:so3-def}
\end{equation}
即\emph{行列式为 $1$ 的正交矩阵}；反之，行列式为 $1$ 的正交矩阵也是旋转矩阵。$9$ 个元素受正交约束 $\mathbf{R}\mathbf{R}^\top=\mathbf{I}$（含 $6$ 个独立标量方程）后只剩 $3$ 个自由度。$\mathrm{SO}$ 即\emph{特殊正交群}（special orthogonal group），“群”的结构留待 \cref{ch:lie}。
\end{definition}

\begin{theorem}[旋转矩阵的正交性]\label{thm:rb-orthonormal}
由 \cref{eq:R-derive} 构造的 $\mathbf{R}$ 满足 $\mathbf{R}^\top\mathbf{R}=\mathbf{I}$ 且 $\det\mathbf{R}=+1$，故 $\mathbf{R}\in\SO(3)$。
\end{theorem}
\begin{proof}
此结论各书或留作习题（如视觉SLAM十四讲\cite{gaoxiang2019slam14}）、或直接陈述，本书补出完整代数证明。$\mathbf{R}$ 的 $(i,j)$ 元为 $R_{ij}=\mathbf{e}_i^\top\mathbf{e}_j'$。计算 $\mathbf{R}^\top\mathbf{R}$ 的 $(j,k)$ 元：
\[
  (\mathbf{R}^\top\mathbf{R})_{jk}=\sum_{i=1}^3 R_{ij}R_{ik}=\sum_{i=1}^3(\mathbf{e}_i^\top\mathbf{e}_j')(\mathbf{e}_i^\top\mathbf{e}_k')
  =\mathbf{e}_j'^\top\Big(\sum_{i=1}^3\mathbf{e}_i\mathbf{e}_i^\top\Big)\mathbf{e}_k'.
\]
因 $\{\mathbf{e}_i\}$ 是一组标准正交\emph{完备}基，有恒等式 $\sum_i\mathbf{e}_i\mathbf{e}_i^\top=\mathbf{I}$（外积之和=单位阵）。代入得 $(\mathbf{R}^\top\mathbf{R})_{jk}=\mathbf{e}_j'^\top\mathbf{e}_k'=\delta_{jk}$，即 $\mathbf{R}^\top\mathbf{R}=\mathbf{I}$，$\mathbf{R}$ 正交。又因 $\{\mathbf{e}_i\}$ 与 $\{\mathbf{e}_i'\}$ 同为右手系（旋转不含镜像），由 $\det(\mathbf{R}^\top\mathbf{R})=(\det\mathbf{R})^2=1$ 得 $\det\mathbf{R}=\pm1$，而右手系到右手系排除 $-1$，故 $\det\mathbf{R}=+1$。Barfoot 给出等价的第二条路线：从轴角式 \cref{eq:rodrigues-rb} 出发直接算 $\mathbf{R}\mathbf{R}^\top$ 并用恒等式 (I3)(I2) 化简为 $\mathbf{I}$（见练习）。
\end{proof}

\begin{proposition}[旋转矩阵的逆是其转置]\label{prop:R-inv}
$\mathbf{R}\in\SO(3)\Rightarrow\mathbf{R}^{-1}=\mathbf{R}^\top\in\SO(3)$，几何上对应“反向旋转”：$\mathbf{a}'=\mathbf{R}^{-1}\mathbf{a}=\mathbf{R}^\top\mathbf{a}$。
\end{proposition}
\begin{proof}
由正交性 $\mathbf{R}\mathbf{R}^\top=\mathbf{I}$ 直接得 $\mathbf{R}^{-1}=\mathbf{R}^\top$。又 $(\mathbf{R}^\top)(\mathbf{R}^\top)^\top=\mathbf{R}^\top\mathbf{R}=\mathbf{I}$ 且 $\det\mathbf{R}^\top=\det\mathbf{R}=1$，故 $\mathbf{R}^\top\in\SO(3)$。在 vectrix 记号下这表现为 $\mathbf{C}_{12}=\mathbf{C}_{21}^{-1}=\mathbf{C}_{21}^\top$\cite{barfoot2024state}。
\end{proof}

\paragraph{复合（链式相消）。}三个系 $1,2,3$，由 $\mathbf{r}_3=\mathbf{R}_{32}\mathbf{r}_2=\mathbf{R}_{32}\mathbf{R}_{21}\mathbf{r}_1$ 且 $\mathbf{r}_3=\mathbf{R}_{31}\mathbf{r}_1$ 得 $\mathbf{R}_{31}=\mathbf{R}_{32}\mathbf{R}_{21}$——\emph{相邻下标相消}，这是后面记号约定一节的记忆诀窍。

\paragraph{主动 vs 被动。}\cref{eq:R-derive} 是\emph{被动}解读（alias：向量不动，坐标系旋转，坐标随之变换）。同一个 $\mathbf{R}$ 也可\emph{主动}解读（alibi：坐标系不动，把点 $\mathbf{a}$ 旋转到新位置 $\mathbf{R}\mathbf{a}$）。两种解读数值相同、语义相反，务必在具体问题中明确取哪种。本书坐标变换取被动解读（$\mathbf{p}_w=\mathbf{R}_{wc}\mathbf{p}_c$），而 \cref{sec:rb-quat} 的四元数旋转点是主动解读。

\paragraph{主轴旋转矩阵（Barfoot 显式给出，本书补全）。}绕单个基轴的旋转是最基本的“积木”，欧拉角即由它们相乘而成。Barfoot 用\emph{被动（frame-rotation）}约定写出三者\cite{barfoot2024state}（$\mathbf{C}_i$ 表绕第 $i$ 轴转 $\theta_i$，把系 $1$ 坐标变到系 $2$）：
\begin{align}
  \mathbf{C}_3(\theta_3)&=\begin{bmatrix}\cos\theta_3&\sin\theta_3&0\\-\sin\theta_3&\cos\theta_3&0\\0&0&1\end{bmatrix},\quad
  \mathbf{C}_2(\theta_2)=\begin{bmatrix}\cos\theta_2&0&-\sin\theta_2\\0&1&0\\\sin\theta_2&0&\cos\theta_2\end{bmatrix},\notag\\
  \mathbf{C}_1(\theta_1)&=\begin{bmatrix}1&0&0\\0&\cos\theta_1&\sin\theta_1\\0&-\sin\theta_1&\cos\theta_1\end{bmatrix}.
  \label{eq:rb-principal}
\end{align}
\textbf{符号布局注意}：这是\emph{被动}约定，$+\sin$ 在“右上”（如 $\mathbf{C}_3$ 的 $(1,2)$ 元）；十四讲常用的\emph{主动}约定 $\mathbf{R}_z(\theta)=[\cos\theta,-\sin\theta,0;\sin\theta,\cos\theta,0;0,0,1]$ 把 $+\sin$ 放“左下”，恰好是转置关系（$\mathbf{R}_z(\theta)=\mathbf{C}_3(\theta)^\top=\mathbf{C}_3(-\theta)$）。这两套差一个转置/取负，混用会让所有旋转方向反过来——务必在一个推导里始终用同一套。

\begin{pitfall}[概念误区：把 $\mathbf{R}^{-1}=\mathbf{R}^\top$ 当巧合]
$\mathbf{R}^{-1}=\mathbf{R}^\top$ 不是“凑巧好算”，而是正交性的\emph{定义}后果：$\mathbf{R}^\top\mathbf{R}=\mathbf{I}$ 等价于“$\mathbf{R}$ 的列向量构成一组标准正交基”，几何上即旋转保长度（保范），故其逆就是“反着转”，而反着转恰是转置。哪些矩阵逆等于转置？——正交矩阵；逆等于共轭转置？——酉矩阵（量子力学中的旋转）。理解了这点，就明白为何下一节 $\SE(3)$ 的逆\emph{不是}简单转置（多了平移项 $-\mathbf{R}^\top\mathbf{t}$，因为 $\SE(3)$ 不是正交群）。
\end{pitfall}

\begin{practice}
\begin{enumerate}
  \item （验证）写出主动约定下绕 $z$ 轴转 $\theta$ 的 $\mathbf{R}_z(\theta)$，验证 $\mathbf{R}_z\mathbf{R}_z^\top=\mathbf{I}$、$\det=1$，并说明它与 $\mathbf{C}_3(\theta)$ 的转置关系。
  \item （证明）证明 $\det\mathbf{R}=\pm1$ 对一切正交阵成立；$\det=-1$ 的正交阵对应什么几何变换？\emph{答案线索}：镜像/反射。
  \item （证明，承重）从轴角式 \cref{eq:rodrigues-rb} 出发，用 (I2)(I3) 证明 $\mathbf{R}\mathbf{R}^\top=\mathbf{I}$（Barfoot 习题 7.2 的正交性第二证）。
\end{enumerate}
\end{practice}

% ════════════════════════════════════════════════════════════════
\section{欧氏变换与齐次坐标}
\label{sec:rb-se3}
\noindent\emph{节首导读：骨干必跟。把“旋转 + 平移”塞进一个矩阵，是相机模型、位姿图、轨迹估计的通用语言。}

\paragraph{加入平移。}世界系向量 $\mathbf{a}$ 经旋转 $\mathbf{R}$ 与平移 $\mathbf{t}$ 得
\begin{equation}
  \mathbf{a}'=\mathbf{R}\mathbf{a}+\mathbf{t}.
  \label{eq:euclid}
\end{equation}
$\mathbf{t}$ 称\emph{平移向量}。用一个 $\mathbf{R}$ 和一个 $\mathbf{t}$ 即可完整描述欧氏空间的坐标变换。写成双系记法（向量 $\mathbf{a}$ 在系 $1,2$ 下坐标 $\mathbf{a}_1,\mathbf{a}_2$）：$\mathbf{a}_1=\mathbf{R}_{12}\mathbf{a}_2+\mathbf{t}_{12}$\cite{gaoxiang2019slam14}。

\paragraph{反面：连续变换会“繁琐”，且非线性。}两次变换 $\mathbf{b}=\mathbf{R}_1\mathbf{a}+\mathbf{t}_1$、$\mathbf{c}=\mathbf{R}_2\mathbf{b}+\mathbf{t}_2$ 串起来是 $\mathbf{c}=\mathbf{R}_2(\mathbf{R}_1\mathbf{a}+\mathbf{t}_1)+\mathbf{t}_2$，变换多次后嵌套很深，且 \cref{eq:euclid} 不是关于 $\mathbf{a}$ 的\emph{线性}映射（有常数项 $\mathbf{t}$）。解决办法：\textbf{齐次坐标}——在三维坐标末尾补 $1$ 变四维，把旋转与平移塞进一个矩阵：
\begin{equation}
  \begin{bmatrix}\mathbf{a}'\\1\end{bmatrix}
  =\underbrace{\begin{bmatrix}\mathbf{R}&\mathbf{t}\\\mathbf{0}^\top&1\end{bmatrix}}_{\mathbf{T}}
  \begin{bmatrix}\mathbf{a}\\1\end{bmatrix}.
  \label{eq:homo}
\end{equation}
于是变换变成线性的，连续变换化为矩阵连乘：用 $\tilde{\mathbf{a}}$ 表 $\mathbf{a}$ 的齐次坐标，$\tilde{\mathbf{b}}=\mathbf{T}_1\tilde{\mathbf{a}}$、$\tilde{\mathbf{c}}=\mathbf{T}_2\tilde{\mathbf{b}}$，故 $\tilde{\mathbf{c}}=\mathbf{T}_2\mathbf{T}_1\tilde{\mathbf{a}}$。约定：不引起歧义时直接写 $\mathbf{b}=\mathbf{T}\mathbf{a}$，默认已做齐次坐标转换（C++ 中可用运算符重载统一齐次/非齐次写法）。

\begin{definition}[特殊欧氏群 $\SE(3)$]\label{def:se3-rep}
\begin{equation}
  \SE(3)=\Big\{\mathbf{T}=\begin{bmatrix}\mathbf{R}&\mathbf{t}\\\mathbf{0}^\top&1\end{bmatrix}\in\mathbb{R}^{4\times4}\,\Big|\,\mathbf{R}\in\SO(3),\ \mathbf{t}\in\mathbb{R}^3\Big\}.
\end{equation}
$16$ 个元素表 $6$ 个自由度（$3$ 旋转 $+$ $3$ 平移）。$\mathrm{SE}$ 即\emph{特殊欧氏群}（special Euclidean group）。
\end{definition}

\begin{proposition}[$\SE(3)$ 的逆]\label{prop:T-inv}
\begin{equation}
  \mathbf{T}^{-1}=\begin{bmatrix}\mathbf{R}^\top&-\mathbf{R}^\top\mathbf{t}\\\mathbf{0}^\top&1\end{bmatrix}.
  \label{eq:T-inv}
\end{equation}
\end{proposition}
\begin{proof}
直接验证（十四讲、Barfoot 均直接给结果，这里补出验证）：
\[
  \mathbf{T}\mathbf{T}^{-1}=\begin{bmatrix}\mathbf{R}&\mathbf{t}\\\mathbf{0}^\top&1\end{bmatrix}\begin{bmatrix}\mathbf{R}^\top&-\mathbf{R}^\top\mathbf{t}\\\mathbf{0}^\top&1\end{bmatrix}
  =\begin{bmatrix}\mathbf{R}\mathbf{R}^\top&-\mathbf{R}\mathbf{R}^\top\mathbf{t}+\mathbf{t}\\\mathbf{0}^\top&1\end{bmatrix}
  =\begin{bmatrix}\mathbf{I}&\mathbf{0}\\\mathbf{0}^\top&1\end{bmatrix}.
\]
注意逆的平移项是 $-\mathbf{R}^\top\mathbf{t}$，\emph{不是} $-\mathbf{t}$——这正是下一节“$\mathbf{t}_{wc}\neq-\mathbf{t}_{cw}$”的根源。Barfoot 在 vectrix 记号下把同一结果写成 $\mathbf{T}_{iv}^{-1}=\mathbf{T}_{vi}$，平移块经 $\mathbf{r}_v^{iv}=-\mathbf{C}_{iv}^\top\mathbf{r}_i^{vi}$ 翻转\cite{barfoot2024state}。
\end{proof}

\paragraph{复合变换与“先旋转后平移”。}与旋转一样，位姿复合即矩阵连乘且相邻下标相消：$\mathbf{T}_{iv}=\mathbf{T}_{ia}\mathbf{T}_{ab}\mathbf{T}_{bv}$。逐块展开 $\mathbf{T}_2\mathbf{T}_1=\big[\mathbf{R}_2\mathbf{R}_1,\ \mathbf{R}_2\mathbf{t}_1+\mathbf{t}_2;\ \mathbf{0}^\top,1\big]$，可见旋转块相乘、平移块是“先把 $\mathbf{t}_1$ 旋到新系再加 $\mathbf{t}_2$”——这就是“变换矩阵先旋转后平移”的约定\cite{barfoot2024state}。

\begin{note}
\textbf{齐次坐标的两个用途.}\;(1) 把仿射变换线性化（本节）。(2) 携带\emph{尺度因子}：齐次点 $[sx,sy,sz,s]^\top$ 与 $[x,y,z,1]^\top$ 表示同一三维点（前三除以第四即恢复），$s\to0$ 可表“无穷远点”——这在相机模型与射影几何里至关重要。齐次坐标由 Möbius（1827）在《重心计算》中引入：把平面点用三角形三顶点的“质量”重心表示，坐标按比例缩放不变，曲线方程在此坐标下变齐次\cite{barfoot2024state}。
\end{note}

\begin{practice}
\begin{enumerate}
  \item 用 \cref{eq:homo} 把 $\mathbf{c}=\mathbf{R}_2(\mathbf{R}_1\mathbf{a}+\mathbf{t}_1)+\mathbf{t}_2$ 写成 $\tilde{\mathbf{c}}=\mathbf{T}_2\mathbf{T}_1\tilde{\mathbf{a}}$，并算出 $\mathbf{T}_2\mathbf{T}_1$ 的旋转块与平移块。
  \item 由 \cref{eq:T-inv} 解释：为什么知道“相机到世界” $\mathbf{T}_{wc}$ 就能立刻得到“世界到相机” $\mathbf{T}_{cw}=\mathbf{T}_{wc}^{-1}$。
  \item （概念）齐次点 $[2,4,6,2]^\top$ 表示哪个三维点？$[1,1,1,0]^\top$ 表示什么？
\end{enumerate}
\end{practice}

\paragraph{阶段回顾与过渡。}至此，本章前三节已经把刚体运动的“静态表示”搭好了一条完整链条：先区分了\emph{向量}与它依赖基的\emph{坐标}（\cref{sec:rb-vector}），再由基变换导出\emph{旋转矩阵} $\SO(3)$ 描述朝向、用 $\mathbb{R}^3$ 向量描述平移（\cref{sec:rb-rotation}），最后把二者塞进一个 $4\times4$ 矩阵得到\emph{位姿} $\SE(3)$，把仿射变换线性化、把连续变换化为矩阵连乘（本节）。表示已立，但有一个隐患还没解决：多坐标系一旦出现，下标 $\mathbf{T}_{wc}$ 究竟读“相机到世界”还是“世界到相机”、平移块到底是谁指向谁——这类记号歧义是 SLAM 实践中 bug 的头号来源。下一节就专门把全书的记号约定一次讲死，并与 Barfoot、Handbook 对照，免得日后读不同文献时把方向写反。

% ════════════════════════════════════════════════════════════════
\section{坐标系与位姿的记号约定}
\label{sec:rb-convention}
\noindent\emph{节首导读：骨干必跟。记号混乱是 SLAM 实践中误解与 bug 的头号来源\cite{barfoot2024state}。本节把全书约定一次讲清并对照三书。}

\paragraph{本书约定：下标从右读到左，左为目标。}
$\mathbf{R}_{wc}$ 表示“把\emph{相机系} $c$ 的坐标变到\emph{世界系} $w$”：
\begin{equation}
  \mathbf{p}_w=\mathbf{R}_{wc}\,\mathbf{p}_c,\qquad \mathbf{p}_w=\mathbf{T}_{wc}\,\mathbf{p}_c\ (\text{齐次}).
\end{equation}
相邻下标相消是记忆复合方向的诀窍：$\mathbf{T}_{wb}=\mathbf{T}_{wc}\mathbf{T}_{cb}$（中间的 $c$ 相消）。本书把坐标系标号写在\emph{右下标}。

\paragraph{平移的语义陷阱。}$\mathbf{t}_{wc}$ 是“\emph{相机系原点}在\emph{世界系}下的坐标”，即相机在世界中的位置——所以 $\mathbf{T}_{wc}$ 直观（平移块直接告诉你相机在哪）。$\mathbf{t}_{12}$ 读作“从 $1$ 到 $2$ 的向量”，即坐标系 $1$ 原点指向坐标系 $2$ 原点的向量、在坐标系 $1$ 下取的坐标\cite{gaoxiang2019slam14}。\textbf{关键陷阱：$\mathbf{t}_{wc}\neq-\mathbf{t}_{cw}$}：由 \cref{eq:T-inv}，$\mathbf{t}_{cw}=-\mathbf{R}_{cw}\mathbf{t}_{wc}=-\mathbf{R}_{wc}^\top\mathbf{t}_{wc}$，两者还差一个旋转。

\begin{pitfall}[编程陷阱：平移直接取负]
新手常把“世界到相机”的平移写成“相机到世界”平移的相反数。错！\textbf{现象}：坐标变换结果位置整体偏到错误的地方，但旋转看起来对，调试时极易误判为“数据问题”。\textbf{根本原因}：$\mathbf{T}^{-1}$ 的平移项是 $-\mathbf{R}^\top\mathbf{t}$ 而非 $-\mathbf{t}$；只有\emph{旋转}部分满足简单转置 $\mathbf{R}_{cw}=\mathbf{R}_{wc}^\top$。\textbf{正确做法}：位姿求逆务必用 \cref{eq:T-inv} 整体来算（或直接调库的 \texttt{inverse()}），绝不要分别对 $\mathbf{R}$、$\mathbf{t}$ 取负。\textbf{自检}：算一遍 $\mathbf{T}\mathbf{T}^{-1}$ 看是否为单位阵。
\end{pitfall}

\paragraph{三书记号对照（查阅原文必备）。}
\begin{table}[H]\centering\small
\caption{坐标系/位姿记号对照：同一对象，三种写法}
\label{tab:rb-frame-notation}
\begin{tabular}{llll}
\toprule
对象 & 本书 & Barfoot\cite{barfoot2024state} & Handbook\cite{carlone2026handbook} \\
\midrule
旋转矩阵 & $\mathbf{R}_{wc}$（左目标右源）& $\mathbf{C}_{wc}$（同左目标右源）& $\mathbf{R}_c^{w}$（\textbf{下源上目标}）\\
作用式 & $\mathbf{p}_w=\mathbf{R}_{wc}\mathbf{p}_c$ & $\mathbf{p}_w=\mathbf{C}_{wc}\mathbf{p}_c$ & $\boldsymbol{\ell}^w=\mathbf{R}_c^{w}\boldsymbol{\ell}^c$ \\
向量坐标的帧名 & 下标 $\mathbf{p}_c$ & 下标 $\mathbf{r}_c$ & \textbf{上标} $\boldsymbol{\ell}^c$ \\
位姿 & $\mathbf{T}_{wc}\in\SE(3)$ & $\mathbf{T}_{wc}$ & $\mathbf{T}_c^{w}$ \\
旋转字母 & $\mathbf{R}$ & $\mathbf{C}$ & $\mathbf{R}$ \\
\bottomrule
\end{tabular}
\end{table}

\noindent 三者几何含义\emph{完全相同}（“目标 $\leftarrow$ 源”的复合、相邻帧相消），仅 Handbook 把目标帧放\emph{上标}、源帧放\emph{下标}，且向量坐标的帧名也写成上标。抄录 Handbook 公式入本书时，把 $\mathbf{R}_a^b\mapsto\mathbf{R}_{ba}$（交换上下标层次）即可，否则\emph{所有坐标变换方向都会反}。

\paragraph{Barfoot 的“左手版”机器人约定（务必识别）。}Barfoot 在“机器人约定”一节里，为了能直接用“车体相对世界”的角 $\theta_{vi}$ 表朝向，把轴角公式第三项写成 $+\sin\theta\,\mathbf{a}^\wedge$（而非主公式的 $-\sin\phi\,\mathbf{a}^\times$），并\emph{明说这对应左手旋转}\cite{barfoot2024state}。读 Barfoot 时若见到 $+\sin$ 版 \cref{eq:rodrigues-rb} 配 $\mathbf{C}_{iv}$，要意识到它与“被动 $\mathbf{C}_{21}$、$-\sin$ 版”差一个转置；本书统一到主动点旋转的 $+\sin$ 版（见 \cref{thm:rodrigues-rb}），引用时按帧对齐即可。

\begin{insight}
为什么本书选 $\mathbf{T}_{wc}$（相机到世界）而非 $\mathbf{T}_{cw}$？因为机器人“轨迹”就是相机系原点在世界中的位置 $\mathbf{O}_w=\mathbf{T}_{wc}\mathbf{O}_c=\mathbf{t}_{wc}$（$\mathbf{O}_c=\mathbf{0}$）——$\mathbf{T}_{wc}$ 的平移块\emph{直接}就是相机在世界中的坐标，可视化、调试都更直观\cite{gaoxiang2019slam14}。这也是轨迹文件常存 $\mathbf{T}_{wc}$ 而非 $\mathbf{T}_{cw}$ 的原因（\cref{sec:rb-eigen} 的可视化例）。
\end{insight}

\begin{practice}
\begin{enumerate}
  \item 已知 $\mathbf{T}_{wc_1},\mathbf{T}_{wc_2}$，写出把点从 $c_1$ 系变到 $c_2$ 系的复合表达式（用相邻相消验证方向）。
  \item Handbook 给你 $\mathbf{R}_a^b$，本书要 $\mathbf{R}_{ab}$，二者哪个把 $a$ 系坐标变到 $b$ 系？哪个反过来？\emph{答案线索}：$\mathbf{R}_a^b=\mathbf{R}_{ba}$。
\end{enumerate}
\end{practice}

% ════════════════════════════════════════════════════════════════
\section{旋转向量与罗德里格斯公式}
\label{sec:rb-axisangle}
\noindent\emph{节首导读：骨干必跟（难度 $\bigstar\bigstar\bigstar$）。轴角是最紧凑的旋转描述，也是 \cref{ch:lie} 李代数的几何原型；Rodrigues 公式是本章的核心定理。}

\paragraph{动机：旋转矩阵太“胖”。}$\SO(3)$ 用 $9$ 个量表 $3$ 自由度，\emph{冗余}且\emph{带约束}（正交、$\det=1$），做存储和（未来的）优化都不便\cite{gaoxiang2019slam14}。最紧凑的几何描述是\textbf{轴角}（旋转向量）：由\emph{欧拉旋转定理}（刚体绕一固定点的最一般运动是绕过该点某轴的一次旋转\cite{barfoot2024state}），任意旋转都可由一个\emph{单位转轴} $\mathbf{a}$ 与一个\emph{转角} $\theta$ 刻画，合写成一个三维向量 $\boldsymbol{\phi}=\theta\mathbf{a}$（方向为轴、模长为角）。变换则用一个旋转向量 $+$ 一个平移向量（共 $6$ 维）。

\begin{theorem}[罗德里格斯公式：轴角 $\to$ 旋转矩阵]\label{thm:rodrigues-rb}
绕单位轴 $\mathbf{a}$（$\mathbf{a}^\top\mathbf{a}=1$）转 $\theta$ 的旋转矩阵为
\begin{equation}
  \mathbf{R}=\cos\theta\,\mathbf{I}+(1-\cos\theta)\,\mathbf{a}\mathbf{a}^\top+\sin\theta\,\mathbf{a}^\wedge.
  \label{eq:rodrigues-rb}
\end{equation}
\end{theorem}
\begin{insight}[两条路殊途同归：几何给“为什么”，代数给“算什么”]
理解 Rodrigues 公式的最好办法，是\emph{同时}走两条互补的路、再看它们在终点汇合——这也是 Craig\cite{craig2005introduction}、\emph{Modern Robotics}\cite{lynch2017modern} 与 Siciliano\cite{siciliano2009robotics} 各自侧重的两个视角：
\begin{itemize}
  \item \textbf{几何动机一路（回答“为什么是这个样子”）}：把待转向量按转轴\emph{分解}为平行分量 $\mathbf{p}_\parallel$ 与垂直分量 $\mathbf{p}_\perp$。绕轴转动\emph{只}动垂直分量——它躺在与轴正交的\emph{切平面}内，在该平面里规规矩矩地转过 $\theta$；平行分量则纹丝不动。于是 $\sin\theta$ 与 $\cos\theta$ 不是从天而降的系数，而是切平面内一个二维转动的横纵投影，$\mathbf{a}^\wedge\mathbf{p}$ 不过是切平面里那条与 $\mathbf{p}_\perp$ 等长且超前 $90^\circ$ 的“虚轴”。这条路把公式的每一项都\emph{画得出来}，却不便机械推广（详见\textbf{证明一}）。
  \item \textbf{代数形式一路（回答“具体怎么算”）}：直接把旋转写成反对称阵的矩阵指数，利用 $\mathbf{a}^\wedge$ 的幂只在 $\{\mathbf{a}^\wedge,(\mathbf{a}^\wedge)^2\}$ 间循环（恒等式 \cref{eq:rb-id3,eq:rb-id4}），把级数按奇偶项归并、认出 $\sin\theta$ 与 $1-\cos\theta$，即得最便于编程的\emph{规范代数形式}
  \[
    \exp(\theta\,\mathbf{a}^\wedge)=\mathbf{I}+\sin\theta\,\mathbf{a}^\wedge+(1-\cos\theta)(\mathbf{a}^\wedge)^2,
  \]
  再用 \cref{eq:rb-id3} 把 $(\mathbf{a}^\wedge)^2=\mathbf{a}\mathbf{a}^\top-\mathbf{I}$ 回代，便与 \cref{eq:rodrigues-rb} 完全一致。这条路\emph{看不见图}，却把“轴角空间正是李代数 $\so(3)$、旋转是它的指数映射”这件深层结构和盘托出，并且无脑可算（详见\textbf{证明二}）。
\end{itemize}
\noindent 两路在 \cref{eq:rodrigues-rb} 处汇合，互为印证：几何法告诉你公式\emph{为什么}长这样、每项的物理含义是什么，代数法告诉你公式\emph{怎么算}、并接通 \cref{ch:lie} 的指数映射框架。记住其一即可应用，两者并看才算真懂。
\end{insight}

\noindent 此式在不少文献中只给结果（如视觉SLAM十四讲留作习题\cite{gaoxiang2019slam14}、Barfoot 仅以欧拉旋转定理\emph{陈述}此式\cite{barfoot2024state}）。为使本书自包含，下面给出\emph{两个方向}的完整证明：几何分解法（与本章外积约定一致）与矩阵指数级数法（呼应 \cref{ch:lie}）。

\begin{proof}[证明一：几何分解法（向量平行/垂直分量）]
把待旋转向量 $\mathbf{p}$ 分解为沿轴的分量与垂直于轴的分量：
\[
  \mathbf{p}_\parallel=(\mathbf{a}^\top\mathbf{p})\mathbf{a}=\mathbf{a}\mathbf{a}^\top\mathbf{p},\qquad
  \mathbf{p}_\perp=\mathbf{p}-\mathbf{a}\mathbf{a}^\top\mathbf{p}=(\mathbf{I}-\mathbf{a}\mathbf{a}^\top)\mathbf{p}.
\]
绕轴旋转时 $\mathbf{p}_\parallel$ 不变；$\mathbf{p}_\perp$ 在垂直平面内转 $\theta$。该平面内取一对正交方向：$\mathbf{p}_\perp$ 本身，以及 $\mathbf{a}\times\mathbf{p}=\mathbf{a}^\wedge\mathbf{p}$（它与 $\mathbf{p}_\perp$ 等长且垂直，因 $|\mathbf{a}^\wedge\mathbf{p}|=|\mathbf{a}||\mathbf{p}|\sin\angle=|\mathbf{p}_\perp|$）。于是
\[
  \mathbf{p}_\perp\ \longmapsto\ \cos\theta\,\mathbf{p}_\perp+\sin\theta\,(\mathbf{a}^\wedge\mathbf{p}).
\]
合并旋转后的向量：
\begin{align*}
  \mathbf{p}'&=\mathbf{p}_\parallel+\cos\theta\,\mathbf{p}_\perp+\sin\theta\,\mathbf{a}^\wedge\mathbf{p}
   =\mathbf{a}\mathbf{a}^\top\mathbf{p}+\cos\theta\,(\mathbf{p}-\mathbf{a}\mathbf{a}^\top\mathbf{p})+\sin\theta\,\mathbf{a}^\wedge\mathbf{p}\\
  &=\big[\cos\theta\,\mathbf{I}+(1-\cos\theta)\,\mathbf{a}\mathbf{a}^\top+\sin\theta\,\mathbf{a}^\wedge\big]\mathbf{p}.
\end{align*}
对任意 $\mathbf{p}$ 成立，即 \cref{eq:rodrigues-rb}。
\end{proof}

\begin{proof}[证明二：矩阵指数级数法（呼应 \cref{ch:lie}）]
绕轴 $\theta\mathbf{a}$ 的旋转可写为矩阵指数 $\mathbf{R}=\exp(\theta\mathbf{a}^\wedge)=\sum_{k\ge0}\frac{(\theta\mathbf{a}^\wedge)^k}{k!}$（这一点的来由——为什么旋转等于反对称阵的指数——在 \cref{ch:lie} 由 $\mathbf{R}\mathbf{R}^\top=\mathbf{I}$ 求导得到；此处先承认它）。用单位向量恒等式 (I3) $\mathbf{a}^\wedge\mathbf{a}^\wedge=\mathbf{a}\mathbf{a}^\top-\mathbf{I}$ 与 (I4) $\mathbf{a}^\wedge\mathbf{a}^\wedge\mathbf{a}^\wedge=-\mathbf{a}^\wedge$，归纳得 $\mathbf{a}^\wedge$ 的幂在 $\{\mathbf{a}^\wedge,(\mathbf{a}^\wedge)^2\}$ 间循环：$(\mathbf{a}^\wedge)^{2m+1}=(-1)^m\mathbf{a}^\wedge$、$(\mathbf{a}^\wedge)^{2m}=(-1)^{m-1}(\mathbf{a}^\wedge)^2\ (m\ge1)$。按奇偶项归并级数：
\begin{align*}
  \exp(\theta\mathbf{a}^\wedge)
  &=\mathbf{I}+\underbrace{\Big(\theta-\tfrac{\theta^3}{3!}+\tfrac{\theta^5}{5!}-\cdots\Big)}_{\sin\theta}\mathbf{a}^\wedge
   +\underbrace{\Big(\tfrac{\theta^2}{2!}-\tfrac{\theta^4}{4!}+\cdots\Big)}_{1-\cos\theta}(\mathbf{a}^\wedge)^2\\
  &=\mathbf{I}+\sin\theta\,\mathbf{a}^\wedge+(1-\cos\theta)(\mathbf{a}\mathbf{a}^\top-\mathbf{I})
   =\cos\theta\,\mathbf{I}+(1-\cos\theta)\mathbf{a}\mathbf{a}^\top+\sin\theta\,\mathbf{a}^\wedge,
\end{align*}
末步用 (I3) 回代，得 \cref{eq:rodrigues-rb}。两证殊途同归：几何法直观，级数法揭示“轴角空间正是 $\SO(3)$ 的李代数 $\so(3)$”。
\end{proof}

\begin{note}
\cref{ch:lie} 会从矩阵指数 $\mathbf{R}=\exp(\boldsymbol{\phi}^\wedge)$ 的角度再次得到同一公式，并把这里的轴角 $\boldsymbol{\phi}=\theta\mathbf{a}$ 直接当作 $\so(3)$ 元素。本节的两个证明正是那里“指数映射”的具体计算。
\end{note}

\paragraph{反向：旋转矩阵 $\to$ 轴角（完整推导）。}对 \cref{eq:rodrigues-rb} 两边取迹，逐项用 (I2) 与 $\operatorname{tr}(\mathbf{a}\mathbf{a}^\top)=\mathbf{a}^\top\mathbf{a}=1$、$\operatorname{tr}(\mathbf{a}^\wedge)=0$：
\begin{equation}
  \operatorname{tr}(\mathbf{R})=\cos\theta\cdot3+(1-\cos\theta)\cdot1+\sin\theta\cdot0=3\cos\theta+1-\cos\theta=1+2\cos\theta,
  \label{eq:rb-tracefull}
\end{equation}
于是转角
\begin{equation}
  \theta=\arccos\frac{\operatorname{tr}(\mathbf{R})-1}{2}.
  \label{eq:rb-trace}
\end{equation}
转轴 $\mathbf{a}$ 是旋转的\emph{不动方向}，满足 $\mathbf{R}\mathbf{a}=\mathbf{a}$，即特征值 $1$ 对应的单位特征向量（求解此方程再归一化即得）。直接验证 $\mathbf{R}\mathbf{a}=\cos\theta\,\mathbf{a}+(1-\cos\theta)\mathbf{a}(\mathbf{a}^\top\mathbf{a})+\sin\theta\,\mathbf{a}^\wedge\mathbf{a}=\cos\theta\,\mathbf{a}+(1-\cos\theta)\mathbf{a}+\mathbf{0}=\mathbf{a}$（用 (I2)）。

\begin{pitfall}[思维陷阱：轴角在 $\theta=0,\pi$ 附近的退化]
$\theta\to0$ 时转轴 $\mathbf{a}$ 不确定（任何轴转 $0$ 都是恒等），$\boldsymbol{\phi}=\theta\mathbf{a}\to\mathbf{0}$ 仍良定义但方向丢失；$\theta=\pi$ 时 $\mathbf{a}$ 与 $-\mathbf{a}$ 给出同一旋转，且 \cref{eq:rb-trace} 求 $\mathbf{a}$ 的常用公式分母趋零。这是“三参数表示必有奇异”（\cref{sec:rb-compare}）的体现。\textbf{现象}：从矩阵反求轴角在大角度处数值跳变。\textbf{对策}：数值实现里 $\theta\to0$ 用泰勒展开避免 $\sin\theta/\theta$ 型 $0/0$；$\theta\to\pi$ 时改用更稳的公式（如从 $\mathbf{R}+\mathbf{I}$ 的列取轴）。
\end{pitfall}

\begin{practice}
\begin{enumerate}
  \item （推导）用 \cref{eq:rodrigues-rb} 验证 $\mathbf{R}\mathbf{a}=\mathbf{a}$（提示：$\mathbf{a}^\top\mathbf{a}=1$，$\mathbf{a}^\wedge\mathbf{a}=\mathbf{0}$）。
  \item （推导，承重，草稿纸完成）独立完成证明二的奇偶项归并，得到 \cref{eq:rodrigues-rb}。
  \item 给定 $\mathbf{R}$，写出从 \cref{eq:rb-trace} 求 $\theta$ 后再求单位转轴 $\mathbf{a}$ 的完整步骤，并讨论 $\theta\to\pi$ 时的数值困难。
\end{enumerate}
\end{practice}

% ════════════════════════════════════════════════════════════════
\section{欧拉角与万向锁}
\label{sec:rb-euler}
\noindent\emph{节首导读：次优（理解奇异即可）。欧拉角直观、给人看，但在估计/优化里几乎从不用——理解“为什么不用”比记公式更重要。}

\paragraph{动机：给人看的旋转。}旋转矩阵与轴角对人都不直观。\textbf{欧拉角}把一次旋转拆成绕三个轴的三次转动，每次绕单轴旋转人很容易想象\cite{gaoxiang2019slam14}。但分解方式众多（XYZ、ZYZ、ZYX、$3$-$1$-$3$ 等），且需区分每次绕\emph{固定轴}还是绕\emph{旋转之后的轴}，定义方法极易混淆。

\paragraph{RPY（ZYX，$1$-$2$-$3$ 姿态序列）。}航空航天常用“偏航--俯仰--滚转”（yaw--pitch--roll）。本书约定刚体前方为 $X$ 轴、右侧为 $Y$ 轴、上方为 $Z$ 轴，按 $Z\!-\!Y\!-\!X$ 绕\emph{动轴}依次旋转\cite{gaoxiang2019slam14}：先绕物体 $Z$ 轴转\emph{偏航} $\psi$，再绕\emph{转动后} $Y$ 轴转\emph{俯仰} $\theta$，最后绕\emph{再转动后} $X$ 轴转\emph{滚转} $\phi$。三个角 $[\phi,\theta,\psi]$ 直观可读。Barfoot 用等价的“$1$-$2$-$3$”记法（绕原始 $1$ 轴 roll、中间 $2$ 轴 pitch、变换后 $3$ 轴 yaw），把对应旋转矩阵显式写出\cite{barfoot2024state}（$s_i=\sin\theta_i,c_i=\cos\theta_i$）：
\begin{equation}
  \mathbf{C}_{21}(\theta_3,\theta_2,\theta_1)=\mathbf{C}_3(\theta_3)\mathbf{C}_2(\theta_2)\mathbf{C}_1(\theta_1)
  =\begin{bmatrix}
  c_2c_3 & c_1s_3+s_1s_2c_3 & s_1s_3-c_1s_2c_3\\
  -c_2s_3 & c_1c_3-s_1s_2s_3 & s_1c_3+c_1s_2s_3\\
  s_2 & -s_1c_2 & c_1c_2
  \end{bmatrix}.
  \label{eq:rb-euler123}
\end{equation}
Euler 本人惯用的是 $3$-$1$-$3$ 序列（绕原始 $3$ 轴 $\psi$、中间 $1$ 轴 $\gamma$、变换后 $3$ 轴 $\theta$，经典力学称 $\psi$ 进动、$\gamma$ 章动、$\theta$ 自转），矩阵为 $\mathbf{C}_3(\theta)\mathbf{C}_1(\gamma)\mathbf{C}_3(\psi)$。这些都只是 \cref{eq:rb-principal} 中主轴矩阵的不同顺序乘积。

\paragraph{反面：万向锁（gimbal lock）。}当俯仰 $\theta_2=\pm90^\circ$（即 $\theta_2=\pi/2$）时，第一次与第三次旋转作用在\emph{同一条轴}上，系统从三次独立旋转退化为两次——\textbf{丢失一个自由度}。这就是奇异性。把 $\theta_2=\pi/2$ 代入 \cref{eq:rb-euler123}（$c_2=0,s_2=1$）：
\begin{equation}
  \mathbf{C}_{21}(\theta_3,\tfrac{\pi}{2},\theta_1)=\begin{bmatrix}
  0&\sin(\theta_1+\theta_3)&-\cos(\theta_1+\theta_3)\\
  0&\cos(\theta_1+\theta_3)&\sin(\theta_1+\theta_3)\\
  1&0&0\end{bmatrix},
\end{equation}
矩阵只依赖 $\theta_1+\theta_3$——$\theta_1,\theta_3$ 关联到\emph{同一}自由度，从矩阵反求二者时无法唯一确定\cite{barfoot2024state}。其数学表现在运动学里看得最清楚（\cref{sec:rb-kinematics}）：欧拉角速率映射的逆含 $\sec\theta_2$，在 $\theta_2\to\pm90^\circ$ 处发散。

\begin{insight}[一般性结论：三参数必有奇异（Stuelpnagel 1964）]
不只是某一种欧拉角——\textbf{任何用 $3$ 个实数表示三维旋转的方案，都不可避免地存在奇异点}（Stuelpnagel\cite{stuelpnagel1964parametrization}、Barfoot\cite{barfoot2024state}）。系统结论：参数多于 $3$ 的表示必带\emph{约束}（如旋转矩阵 $9$ 参带正交约束、四元数 $4$ 参带单位约束）；恰好 $3$ 参的表示必带\emph{奇异}（欧拉角、轴角、Gibbs 向量）；\emph{不存在}既最小（$3$ 参）又处处无奇异的完美表示。这条结论决定了 SLAM 里为何用四元数/旋转矩阵存姿态、而非欧拉角。直觉类比：用经纬度两坐标表地球表面必在两极奇异（纬度 $\pm90^\circ$ 时经度无意义）——这个类比在“最小参数化的拓扑障碍”这一点上贴切，但在“维数”上不可延伸（球面是 $2$ 维、$\SO(3)$ 是 $3$ 维）。
\end{insight}

\begin{pitfall}[思维陷阱：在优化/滤波里用欧拉角]
因万向锁与奇异性，欧拉角\textbf{不适合插值、迭代、滤波、优化}。\textbf{现象}：滤波/BA 在大俯仰角处雅可比爆炸、协方差发散。\textbf{根本原因}：速率映射 $\mathbf{S}^{-1}$ 含 $\sec\theta_2$（\cref{sec:rb-kinematics}）。\textbf{正确做法}：SLAM 程序内部用四元数/旋转矩阵表姿态，欧拉角\emph{只}用于人机交互、结果验证（转成欧拉角能快速肉眼判断结果是否合理），或平面场景只取偏航角作定位输出（扫地机、自动驾驶车）。需要时再从四元数/旋转矩阵\emph{转换}出欧拉角给人看。
\end{pitfall}

\begin{practice}
\begin{enumerate}
  \item 把 $\theta_2=\pi/2$ 代入 \cref{eq:rb-euler123}，自行验证矩阵只依赖 $\theta_1+\theta_3$。
  \item （概念）为什么“参数恰好 $3$ 个”与“存在奇异”是同一件事的两面？\emph{答案线索}：$\SO(3)$ 与 $\mathbb{R}^3$ 不同胚。
\end{enumerate}
\end{practice}

% ════════════════════════════════════════════════════════════════
\section{四元数}
\label{sec:rb-quat}
\noindent\emph{节首导读：骨干必跟（难度 $\bigstar\bigstar\bigstar$）。四元数是姿态\emph{存储与插值}的工业首选：紧凑、无奇异。本节全量给出 Hamilton 运算、旋转作用、与矩阵/轴角互转的完整推导。}

\paragraph{动机：既紧凑又无奇异。}既然三参数必有奇异、九参数太冗余，能否用\emph{四个}数无奇异地表旋转？能——\textbf{四元数}\cite{gaoxiang2019slam14}。类比（标边界）：二维旋转可用单位复数 $e^{\mathrm{i}\theta}=\cos\theta+\mathrm{i}\sin\theta$ 表示，复数乘法实现复平面旋转（乘 $\mathrm{i}$ 转 $90^\circ$）；三维旋转可用单位四元数表示，它紧凑、无奇异。类比的\emph{边界}：复数乘法可交换、乘 $\mathrm{i}$ 转 $90^\circ$，而四元数乘法\emph{不可交换}、乘 $\mathrm{i}$ 对应转 $180^\circ$（见双重覆盖洞见），不能把复数的全部直觉照搬过来。

\begin{definition}[四元数（Hamilton 约定）]\label{def:quat}
四元数有一个实部、三个虚部，本书\emph{实部在前}：
\begin{equation}
  \mathbf{q}=q_0+q_1\mathrm{i}+q_2\mathrm{j}+q_3\mathrm{k}=[w,\mathbf{v}]^\top,\quad w=q_0\in\mathbb{R},\ \mathbf{v}=[q_1,q_2,q_3]^\top\in\mathbb{R}^3,
\end{equation}
虚单位满足 \textbf{Hamilton 约定}：
\begin{equation}
  \mathrm{i}^2=\mathrm{j}^2=\mathrm{k}^2=-1,\quad \mathrm{ij}=\mathrm{k},\ \mathrm{ji}=-\mathrm{k},\ \mathrm{jk}=\mathrm{i},\ \mathrm{kj}=-\mathrm{i},\ \mathrm{ki}=\mathrm{j},\ \mathrm{ik}=-\mathrm{j}.
  \label{eq:rb-quat-rules}
\end{equation}
虚部为 $\mathbf{0}$ 称\emph{实四元数}；实部为 $0$ 称\emph{虚四元数}。把 $\mathrm{i,j,k}$ 看成三个坐标轴，它们与自己的乘法和复数一样（平方为 $-1$），相互之间的乘法和\emph{外积}一样（$\mathrm{ij}=\mathrm{k}$ 对应 $\mathbf{e}_1\times\mathbf{e}_2=\mathbf{e}_3$）。
\end{definition}

\begin{note}
\textbf{Hamilton vs JPL（务必分清）.}\;本书取 Hamilton（$\mathrm{ij}=\mathrm{k}$，与 Eigen、十四讲、Barfoot 一致）。航天界常用的 \emph{JPL} 约定取 $\mathrm{ij}=-\mathrm{k}$、实部写在末位，二者旋转复合\emph{顺序相反}，混用必出错。\textbf{排序差异}：本书与十四讲\emph{实部在前} $[w,\mathbf{v}]$；Barfoot 虽同为 Hamilton 惯例，却把\emph{矢部在前}写成 $\mathbf{q}=[\boldsymbol{\varepsilon};\eta]$（$\boldsymbol{\varepsilon}$ 矢、$\eta$ 标），并自有 $4\times4$ 矩阵代数，引用其公式时须重排分量。\textbf{库差异}：Eigen 的 \texttt{Quaterniond} 构造按 $(w,x,y,z)$、但内存 \texttt{coeffs()} 按 $(x,y,z,w)$ 存——书写顺序与存储顺序不同，初学者易混。
\end{note}

\paragraph{Hamilton 运算全套（十四讲六类，全录）。}设 $\mathbf{q}_a=[s_a,\mathbf{v}_a],\mathbf{q}_b=[s_b,\mathbf{v}_b]$。
\begin{itemize}
\item \textbf{加减}（按分量）：$\mathbf{q}_a\pm\mathbf{q}_b=[s_a\pm s_b,\ \mathbf{v}_a\pm\mathbf{v}_b]$。
\item \textbf{乘法}：把每项相乘、虚部按 \cref{eq:rb-quat-rules} 归并，得分量展开
\begin{align}
  \mathbf{q}_a\mathbf{q}_b=&\ (s_as_b-x_ax_b-y_ay_b-z_az_b)\notag\\
  &+(s_ax_b+x_as_b+y_az_b-z_ay_b)\mathrm{i}+(s_ay_b-x_az_b+y_as_b+z_ax_b)\mathrm{j}\notag\\
  &+(s_az_b+x_ay_b-y_ax_b+z_as_b)\mathrm{k},
\end{align}
更简洁的向量形式（含内、外积）：
\begin{equation}
  \mathbf{q}_a\mathbf{q}_b=\big[\,s_as_b-\mathbf{v}_a^\top\mathbf{v}_b,\ \ s_a\mathbf{v}_b+s_b\mathbf{v}_a+\mathbf{v}_a^\wedge\mathbf{v}_b\,\big].
  \label{eq:quat-mul}
\end{equation}
因末项外积，\emph{四元数乘法一般不可交换}（除非 $\mathbf{v}_a,\mathbf{v}_b$ 共线，外积为零）。两个实四元数乘积仍为实（与复数一致）。
\item \textbf{模长}：$\|\mathbf{q}_a\|=\sqrt{s_a^2+x_a^2+y_a^2+z_a^2}$，且 $\|\mathbf{q}_a\mathbf{q}_b\|=\|\mathbf{q}_a\|\|\mathbf{q}_b\|$（单位四元数相乘仍是单位四元数）。
\item \textbf{共轭}：$\mathbf{q}^*=[s,-\mathbf{v}]$（虚部取反），$\mathbf{q}^*\mathbf{q}=\mathbf{q}\mathbf{q}^*=[s^2+\mathbf{v}^\top\mathbf{v},\,\mathbf{0}]$（实四元数，实部为模长平方）。
\item \textbf{逆}：$\mathbf{q}^{-1}=\mathbf{q}^*/\|\mathbf{q}\|^2$，满足 $\mathbf{q}\mathbf{q}^{-1}=\mathbf{q}^{-1}\mathbf{q}=\mathbf{1}$，且 $(\mathbf{q}_a\mathbf{q}_b)^{-1}=\mathbf{q}_b^{-1}\mathbf{q}_a^{-1}$。\emph{单位四元数的逆等于共轭}。
\item \textbf{数乘}：$k\mathbf{q}=[ks,\,k\mathbf{v}]$。
\end{itemize}

\paragraph{用四元数旋转点。}把点 $\mathbf{p}\in\mathbb{R}^3$ 写成虚四元数 $\mathbf{p}=[0,\mathbf{p}]$（$3$ 个虚部对应空间 $3$ 个轴），单位四元数 $\mathbf{q}$ 表示的旋转作用为
\begin{equation}
  \mathbf{p}'=\mathbf{q}\,\mathbf{p}\,\mathbf{q}^{-1}.
  \label{eq:quat-rot}
\end{equation}
可验证结果实部恒为 $0$，故仍是虚四元数，其虚部即旋转后的点（练习）。

\paragraph{四元数乘法的矩阵表示（左乘/右乘）。}定义 $4\times4$ 左乘矩阵 $\mathbf{q}^{+}$ 与右乘矩阵 $\mathbf{q}^{\oplus}$\cite{gaoxiang2019slam14}：
\begin{equation}
  \mathbf{q}^{+}=\begin{bmatrix}s&-\mathbf{v}^\top\\\mathbf{v}&s\mathbf{I}+\mathbf{v}^\wedge\end{bmatrix},\qquad
  \mathbf{q}^{\oplus}=\begin{bmatrix}s&-\mathbf{v}^\top\\\mathbf{v}&s\mathbf{I}-\mathbf{v}^\wedge\end{bmatrix},
  \label{eq:rb-quat-lr}
\end{equation}
二者仅 $\mathbf{v}^\wedge$ 前符号不同。直接验证 $\mathbf{q}_1^{+}\mathbf{q}_2$ 恰好复现 \cref{eq:quat-mul}，故
\begin{equation}
  \mathbf{q}_1\mathbf{q}_2=\mathbf{q}_1^{+}\mathbf{q}_2=\mathbf{q}_2^{\oplus}\mathbf{q}_1.
  \label{eq:rb-quat-lr-eq}
\end{equation}
\begin{note}
\textbf{Barfoot 的等价代数（并列对照）.}\;Barfoot 用矢部在前的 $\mathbf{q}=[\boldsymbol{\varepsilon};\eta]$，对应左/右乘算子写成 $\mathbf{q}^{+}=[\eta\mathbf{I}-\boldsymbol{\varepsilon}^\wedge,\ \boldsymbol{\varepsilon};\ -\boldsymbol{\varepsilon}^\top,\ \eta]$、$\mathbf{q}^{\oplus}=[\eta\mathbf{I}+\boldsymbol{\varepsilon}^\wedge,\ \boldsymbol{\varepsilon};\ -\boldsymbol{\varepsilon}^\top,\ \eta]$，并满足同型恒等式 $\mathbf{u}^{+}\mathbf{v}=\mathbf{v}^{\oplus}\mathbf{u}$、$(\mathbf{u}^{+})^\top=(\mathbf{u}^{+})^{-1}=(\mathbf{u}^{-1})^{+}$、$\mathbf{u}^{+}\mathbf{v}^{\oplus}=\mathbf{v}^{\oplus}\mathbf{u}^{+}$ 等；单位元 $[\mathbf{0};1]$ 满足 $\boldsymbol{\iota}^{+}=\boldsymbol{\iota}^{\oplus}=\mathbf{I}$。两套代数因排序不同而分块位置不同，但描述同一 Hamilton 乘法\cite{barfoot2024state}。
\end{note}

\begin{theorem}[四元数 $\to$ 旋转矩阵]\label{thm:quat2R}
单位四元数 $\mathbf{q}=[s,\mathbf{v}]$（$s^2+\mathbf{v}^\top\mathbf{v}=1$）对应的旋转矩阵为
\begin{equation}
  \mathbf{R}=\mathbf{v}\mathbf{v}^\top+s^2\mathbf{I}+2s\,\mathbf{v}^\wedge+(\mathbf{v}^\wedge)^2.
  \label{eq:quat2R}
\end{equation}
\end{theorem}
\begin{proof}
把 \cref{eq:quat-rot} 用左/右乘矩阵改写（$\mathbf{p}$ 视为四元数）：
\[
  \mathbf{p}'=\mathbf{q}\,\mathbf{p}\,\mathbf{q}^{-1}=\mathbf{q}^{+}\mathbf{p}^{+}\mathbf{q}^{-1}=\mathbf{q}^{+}(\mathbf{q}^{-1})^{\oplus}\mathbf{p},
\]
第二步把 $\mathbf{q}\mathbf{p}$ 写成 $\mathbf{q}^{+}\mathbf{p}$；末步用 \cref{eq:rb-quat-lr-eq} 把右乘 $\mathbf{q}^{-1}$ 写成 $(\mathbf{q}^{-1})^{\oplus}$ 作用在 $\mathbf{p}$ 上。对单位四元数 $\mathbf{q}^{-1}=\mathbf{q}^*=[s,-\mathbf{v}]$，其右乘矩阵 $(\mathbf{q}^{-1})^{\oplus}=[s,\ \mathbf{v}^\top;\ -\mathbf{v},\ s\mathbf{I}+\mathbf{v}^\wedge]$。乘出：
\[
  \mathbf{q}^{+}(\mathbf{q}^{-1})^{\oplus}
  =\begin{bmatrix}s&-\mathbf{v}^\top\\\mathbf{v}&s\mathbf{I}+\mathbf{v}^\wedge\end{bmatrix}
  \begin{bmatrix}s&\mathbf{v}^\top\\-\mathbf{v}&s\mathbf{I}+\mathbf{v}^\wedge\end{bmatrix}
  =\begin{bmatrix}1&\mathbf{0}^\top\\\mathbf{0}&\mathbf{v}\mathbf{v}^\top+s^2\mathbf{I}+2s\mathbf{v}^\wedge+(\mathbf{v}^\wedge)^2\end{bmatrix},
\]
左上角 $s^2+\mathbf{v}^\top\mathbf{v}=1$（单位条件），右上/左下块借单位条件与 $\mathbf{v}^\top\mathbf{v}^\wedge=\mathbf{0}^\top$ 消为零。因 $\mathbf{p},\mathbf{p}'$ 都是虚四元数，右下 $3\times3$ 块即所求 $\mathbf{R}$，得 \cref{eq:quat2R}。展开 $(\mathbf{v}^\wedge)^2$ 用 (I3) 还可写成 Euler 参数形式 $\mathbf{R}=(s^2-\mathbf{v}^\top\mathbf{v})\mathbf{I}+2\mathbf{v}\mathbf{v}^\top+2s\mathbf{v}^\wedge$，与 Barfoot 式一致\cite{barfoot2024state}。
\end{proof}

\paragraph{四元数 $\to$ 轴角（完整推导）。}对 \cref{eq:quat2R} 取迹，逐项：$\operatorname{tr}(\mathbf{v}\mathbf{v}^\top)=\mathbf{v}^\top\mathbf{v}$、$\operatorname{tr}(s^2\mathbf{I})=3s^2$、$\operatorname{tr}(2s\mathbf{v}^\wedge)=0$、$\operatorname{tr}((\mathbf{v}^\wedge)^2)=-2\mathbf{v}^\top\mathbf{v}$（由 (I3)）：
\begin{equation}
  \operatorname{tr}(\mathbf{R})=\mathbf{v}^\top\mathbf{v}+3s^2-2\mathbf{v}^\top\mathbf{v}=3s^2-\mathbf{v}^\top\mathbf{v}\overset{s^2+\mathbf{v}^\top\mathbf{v}=1}{=}3s^2-(1-s^2)=4s^2-1.
\end{equation}
代入 \cref{eq:rb-trace} 得 $\cos\theta=\frac{(4s^2-1)-1}{2}=2s^2-1=2\cos^2\tfrac{\theta}{2}-1$（倍角公式），故 $s=\cos\tfrac{\theta}{2}$，于是
\begin{equation}
  \theta=2\arccos s,\qquad \mathbf{a}=\mathbf{v}\big/\sin\tfrac{\theta}{2}.
  \label{eq:quat2axis}
\end{equation}
即 $\mathbf{q}=[\cos\tfrac{\theta}{2},\ \mathbf{a}\sin\tfrac{\theta}{2}]$。\textbf{反向（轴角 $\to$ 四元数）}即此式倒用——这也正是 Barfoot 的 \emph{Euler 参数} $\eta=\cos\tfrac{\phi}{2},\ \boldsymbol{\varepsilon}=\mathbf{a}\sin\tfrac{\phi}{2}$（堆成 $[\boldsymbol{\varepsilon};\eta]$ 即单位四元数，满足 $\eta^2+\boldsymbol{\varepsilon}^\top\boldsymbol{\varepsilon}=1$）\cite{barfoot2024state}。各表示间的转换在实践中由库直接提供，选最方便的形式即可。

\begin{insight}[双重覆盖：$\mathbf{q}$ 与 $-\mathbf{q}$ 是同一旋转]
由 \cref{eq:quat2axis}，$\theta$ 与 $\theta+2\pi$（即 $\mathbf{q}\to-\mathbf{q}$）给出同一 $\mathbf{R}$：单位四元数到 $\SO(3)$ 是\emph{二对一}（双重覆盖）。这解释了那句玄妙直觉——“乘 $\mathrm{i}$ 对应转 $180^\circ$、$\mathrm{i}^2=-1$ 对应转 $360^\circ$ 得到\emph{相反}的四元数，要转两周才复原”\cite{gaoxiang2019slam14}。这并非四元数的人为缺陷，而是 $\SO(3)$ 自身拓扑的反映：转 $360^\circ$ 与不转，在 $\SO(3)$ 里\emph{不能}连续地相互形变，转 $720^\circ$ 才能——经典的“皮带把戏 / 盘子把戏”（手托盘子转两周手臂不打结）正演示这一点，而单位四元数所在的三维球面 $S^3$ 恰是 $\SO(3)$ 的那个“无扭结”的二重覆盖空间\cite{hanson2006visualizing}。工程上：比较两个姿态四元数前要统一符号（如令 $w\ge0$），否则插值/求差会走“远路”。
\end{insight}

\begin{pitfall}[编程陷阱：四元数未归一化 / 库约定不一致]
\textbf{错误做法}：直接用未归一化的四元数旋转向量，或把 Eigen 的 \texttt{coeffs()} 顺序 $(x,y,z,w)$ 当成 $(w,x,y,z)$ 喂给别的库。\textbf{现象}：旋转结果带尺度/变形，或出现 \texttt{NaN}，或旋转方向莫名其妙错。\textbf{根本原因}：浮点累积使 $\|\mathbf{q}\|\neq1$；不同库的分量顺序、Hamilton/JPL 约定不一致。\textbf{正确做法}：使用前 \texttt{q.\allowbreak normalize()}；跨库传递时显式核对分量顺序与手性约定。\textbf{自检}：打印 $\|\mathbf{q}\|$ 应为 $1$；用同一旋转分别走矩阵与四元数两条路旋转同一向量，结果应一致。
\end{pitfall}

\begin{practice}
\begin{enumerate}
  \item 验证 \cref{eq:quat-rot} 的结果实部为零（提示：用 \cref{eq:quat-mul} 与单位条件）。
  \item （推导）对 \cref{eq:quat2R} 取迹，独立得出 $\operatorname{tr}(\mathbf{R})=4s^2-1$，从而推出 \cref{eq:quat2axis}。
  \item 画一张表，总结旋转矩阵、轴角、欧拉角、四元数四者的两两转换公式（出处指向本章各式）。
\end{enumerate}
\end{practice}

% ════════════════════════════════════════════════════════════════
\section{（选读）Gibbs 向量与各表示总览}
\label{sec:rb-gibbs}
\noindent\emph{节首导读：选读。Gibbs 向量（Cayley--Rodrigues 参数）是又一种三参数表示，Barfoot 用它给出 Rodrigues 公式的两个代数等价性证明，数值上无三角函数。}

\paragraph{定义。}用轴角定义 \textbf{Gibbs 向量}（又名 Cayley--Rodrigues 参数）：
\begin{equation}
  \mathbf{g}=\mathbf{a}\tan\frac{\theta}{2}=\boldsymbol{\varepsilon}/\eta,
  \label{eq:rb-gibbs}
\end{equation}
在 $\theta=\pi$ 处奇异（三参数表示必有奇异）。其旋转矩阵有两种等价形式\cite{barfoot2024state}：
\begin{equation}
  \mathbf{R}=(\mathbf{I}+\mathbf{g}^\wedge)^{-1}(\mathbf{I}-\mathbf{g}^\wedge)
  =\frac{1}{1+\mathbf{g}^\top\mathbf{g}}\big((1-\mathbf{g}^\top\mathbf{g})\mathbf{I}+2\mathbf{g}\mathbf{g}^\top-2\mathbf{g}^\wedge\big).
  \label{eq:rb-cayley}
\end{equation}
注意符号：这里用 Barfoot 的被动（$-\sin$）主公式，第三项为 $-2\mathbf{g}^\wedge$；统一到本书主动约定时第三项变号。

\begin{derivation}[Rodrigues 两向之一：Gibbs 右式 $\to$ 轴角三角式]
代入 $\mathbf{g}=\mathbf{a}\tan(\theta/2)$、$\mathbf{a}^\top\mathbf{a}=1$，\cref{eq:rb-cayley} 右式变为
\[
  \mathbf{R}=\frac{1}{1+\tan^2\tfrac{\theta}{2}}\Big((1-\tan^2\tfrac{\theta}{2})\mathbf{I}+2\tan^2\tfrac{\theta}{2}\,\mathbf{a}\mathbf{a}^\top-2\tan\tfrac{\theta}{2}\,\mathbf{a}^\wedge\Big),
\]
用 $(1+\tan^2\tfrac{\theta}{2})^{-1}=\cos^2\tfrac{\theta}{2}$，并把各系数化为半角恒等式：
\[
  \mathbf{R}=\underbrace{(\cos^2\tfrac{\theta}{2}-\sin^2\tfrac{\theta}{2})}_{\cos\theta}\mathbf{I}
  +\underbrace{2\sin^2\tfrac{\theta}{2}}_{1-\cos\theta}\mathbf{a}\mathbf{a}^\top
  -\underbrace{2\sin\tfrac{\theta}{2}\cos\tfrac{\theta}{2}}_{\sin\theta}\mathbf{a}^\wedge
  =\cos\theta\,\mathbf{I}+(1-\cos\theta)\mathbf{a}\mathbf{a}^\top-\sin\theta\,\mathbf{a}^\wedge,
\]
即被动版轴角式（取转置即本书 \cref{eq:rodrigues-rb}）。
\end{derivation}

\begin{derivation}[Rodrigues 两向之二：Gibbs 两式互证]
用 Neumann 级数 $(\mathbf{I}+\mathbf{g}^\wedge)^{-1}=\sum_{n\ge0}(-\mathbf{g}^\wedge)^n$。关键中间步反复用化简
\[
  (\mathbf{g}^\top\mathbf{g})\mathbf{g}^\wedge=(-\mathbf{g}^\wedge\mathbf{g}^\wedge+\mathbf{g}\mathbf{g}^\top)\mathbf{g}^\wedge
  =-\mathbf{g}^\wedge\mathbf{g}^\wedge\mathbf{g}^\wedge+\mathbf{g}\underbrace{\mathbf{g}^\top\mathbf{g}^\wedge}_{\mathbf{0}^\top}=-\mathbf{g}^\wedge\mathbf{g}^\wedge\mathbf{g}^\wedge,
\]
（用一般向量版 (I3) $\mathbf{g}^\wedge\mathbf{g}^\wedge=-\mathbf{g}^\top\mathbf{g}\,\mathbf{I}+\mathbf{g}\mathbf{g}^\top$ 与 $\mathbf{g}^\top\mathbf{g}^\wedge=\mathbf{0}^\top$）。逐项配对可证
\[
  (1+\mathbf{g}^\top\mathbf{g})(\mathbf{I}+\mathbf{g}^\wedge)^{-1}=\mathbf{I}+\mathbf{g}\mathbf{g}^\top-\mathbf{g}^\wedge,
\]
两边右乘 $(\mathbf{I}-\mathbf{g}^\wedge)$ 并用 (I3) 化 $\mathbf{g}^\wedge\mathbf{g}^\wedge$：
\[
  (1+\mathbf{g}^\top\mathbf{g})\underbrace{(\mathbf{I}+\mathbf{g}^\wedge)^{-1}(\mathbf{I}-\mathbf{g}^\wedge)}_{\mathbf{R}}
  =(\mathbf{I}+\mathbf{g}\mathbf{g}^\top-\mathbf{g}^\wedge)(\mathbf{I}-\mathbf{g}^\wedge)
  =(1-\mathbf{g}^\top\mathbf{g})\mathbf{I}+2\mathbf{g}\mathbf{g}^\top-2\mathbf{g}^\wedge,
\]
两边除以 $(1+\mathbf{g}^\top\mathbf{g})$ 即得 \cref{eq:rb-cayley} 右式，证毕两式相等。
\end{derivation}

\paragraph{各表示总览（含 Euler 参数与 Gibbs）。}
\begin{table}[H]\centering\small
\caption{六种旋转表示与到旋转矩阵的公式（Barfoot 全景，统一到 $\theta,\mathbf{a}$）}
\label{tab:rb-rep-overview}
\begin{tabular}{lcll}
\toprule
表示 & 参数 & 约束/奇异 & $\to$ 旋转矩阵 \\
\midrule
旋转矩阵 & $9$ & 正交+$\det1$（$6$ 约束）；无奇异 & --- \\
Euler 角 $1$-$2$-$3$ & $3$ & $\theta_2=\pi/2$ 奇异 & $\mathbf{C}_3\mathbf{C}_2\mathbf{C}_1$（\cref{eq:rb-euler123}）\\
轴角 $\{\mathbf{a},\theta\}$ & $4\!\to\!3$ & $\theta=0$ 轴未定 & $\cos\theta\mathbf{I}+(1-\cos\theta)\mathbf{a}\mathbf{a}^\top+\sin\theta\mathbf{a}^\wedge$ \\
Euler 参数/四元数 & $4$ & $\|\mathbf{q}\|=1$；无奇异 & $(s^2-\mathbf{v}^\top\mathbf{v})\mathbf{I}+2\mathbf{v}\mathbf{v}^\top+2s\mathbf{v}^\wedge$ \\
Gibbs 向量 $\mathbf{g}$ & $3$ & $\theta=\pi$ 奇异 & $(\mathbf{I}+\mathbf{g}^\wedge)^{-1}(\mathbf{I}-\mathbf{g}^\wedge)$ \\
\bottomrule
\end{tabular}
\end{table}
\noindent 互转要点：轴角 $\to$ 四元数 $\eta=\cos\tfrac{\theta}{2},\boldsymbol{\varepsilon}=\mathbf{a}\sin\tfrac{\theta}{2}$；四元数 $\to$ Gibbs $\mathbf{g}=\boldsymbol{\varepsilon}/\eta$。Gibbs/Euler 参数的最大数值优点：算旋转矩阵\emph{不含三角函数}（航天常用）。

% ════════════════════════════════════════════════════════════════
\section{各旋转表示的对比与选择}
\label{sec:rb-compare}
\noindent\emph{节首导读：次优。把前几节横向拉通，给出工程选型；这条选型链贯穿全书。}

\begin{table}[H]\centering\small
\caption{四种常用旋转表示对比}
\label{tab:rb-compare}
\begin{tabular}{lcccl}
\toprule
表示 & 参数 & 自由度 & 约束/奇异 & 适用 \\
\midrule
旋转矩阵 $\SO(3)$ & $9$ & $3$ & 正交+$\det1$；\emph{无奇异} & 坐标变换、作用于点 \\
轴角 $\boldsymbol{\phi}=\theta\mathbf{a}$ & $3$ & $3$ & 无约束；\emph{有奇异}($\theta=0,\pi$) & 紧凑表达、李代数（\cref{ch:lie}）\\
欧拉角 RPY & $3$ & $3$ & 无约束；\emph{万向锁} & 人机交互、验证、$2$D 偏航 \\
四元数 & $4$ & $3$ & $\|\mathbf{q}\|=1$；\emph{无奇异} & 姿态存储/插值首选 \\
\bottomrule
\end{tabular}
\end{table}

\noindent\textbf{选型经验（系统分类）}：按“无奇异/紧凑/可线性化”三维度取舍——
\begin{itemize}
  \item \emph{存储/传递姿态}用\textbf{四元数}（$4$ 参，紧凑、无奇异、插值平滑）；
  \item \emph{坐标变换、作用于点}用\textbf{旋转矩阵}（直接矩阵乘法，无奇异）；
  \item \emph{要做估计/优化的增量更新}转到\textbf{李代数}（轴角的近亲，\cref{ch:lie}，本书右扰动主线）；
  \item \emph{只在给人看时}才转\textbf{欧拉角}。
\end{itemize}
因不同领域（机器人/视觉/航天）记号与参数化众多、实践中极易误解，本书坚持\emph{单一}一致的记号，引用异制处显式给出换算（与 Barfoot 同此取向\cite{barfoot2024state}）。

% ════════════════════════════════════════════════════════════════
\section{旋转与位姿运动学}
\label{sec:rb-kinematics}
\noindent\emph{节首导读：次优（难度 $\bigstar\bigstar\bigstar$）。表示也会随时间变；本节给出 Poisson 方程、哥氏加速度、Frenet--Serret、独轮车模型，是 IMU 积分与运动模型的基础，并承接 \cref{ch:lie}。}

\paragraph{动机：表示也会随时间变。}机器人在动，$\mathbf{R}(t)$、$\mathbf{T}(t)$ 是时间的函数。它们的导数与\emph{角速度}的关系，是 IMU 积分、运动模型的基础。系 $\underline{\mathcal{F}}_2$ 相对 $\underline{\mathcal{F}}_1$ 转动，角速度记 $\boldsymbol{\omega}_{21}$（$\boldsymbol{\omega}_{12}=-\boldsymbol{\omega}_{21}$），其大小是转速、方向是瞬时转轴\cite{barfoot2024state}。

\paragraph{Poisson 方程（旋转运动学）。}设机体角速度为 $\boldsymbol{\omega}$，则旋转的运动学为
\begin{equation}
  \dot{\mathbf{R}}=\boldsymbol{\omega}^\wedge\mathbf{R}\quad(\boldsymbol{\omega}\text{ 在世界系})\qquad\text{或}\qquad
  \dot{\mathbf{R}}=\mathbf{R}\,\boldsymbol{\omega}^\wedge\quad(\boldsymbol{\omega}\text{ 在机体系}).
  \label{eq:poisson}
\end{equation}
完整推导留在 \cref{ch:lie}（那里由约束 $\mathbf{R}\mathbf{R}^\top=\mathbf{I}$ 求导，得到 $\dot{\mathbf{R}}\mathbf{R}^\top$ 反对称，从而引出李代数）；而\emph{如何保群地数值积分这条方程、如何线性化它、其状态转移矩阵是什么}，则系统地见 \cref{sec:lie-kinematics}。Barfoot 用 vectrix 给出被动型的等价推导\cite{barfoot2024state}：由 $\underline{\mathcal{F}}_1^\top=\underline{\mathcal{F}}_2^\top\mathbf{C}_{21}$ 在惯性系求导、代入基向量随系转动关系 $\dot{\underline{\mathcal{F}}}_2^\top=\boldsymbol{\omega}_{21}\times\underline{\mathcal{F}}_2^\top$，得
\begin{equation}
  \dot{\mathbf{C}}_{21}=-\boldsymbol{\omega}_2^{21\wedge}\mathbf{C}_{21},\qquad
  \boldsymbol{\omega}_2^{21\wedge}=-\dot{\mathbf{C}}_{21}\mathbf{C}_{21}^\top.
  \label{eq:rb-poisson-barfoot}
\end{equation}
Handbook 写成 $\dot{\mathbf{R}}_a^b=\boldsymbol{\omega}^{b\wedge}\mathbf{R}_a^b$\cite{carlone2026handbook}。三者符号差异\emph{完全}来自“$\mathbf{R}$ 方向 + $\boldsymbol{\omega}$ 表达帧 + 左/右乘”三要素的组合，使用时须三者一并对齐。给定测得的 $\boldsymbol{\omega}$ 即可积分得 $\mathbf{R}(t)$——这就是\emph{捷联导航}（strapdown navigation），因测角速度的陀螺固连在转动系上。

\paragraph{哥氏关系与加速度（Barfoot 完整给出）。}设 $(\cdot)^{\bullet}$ 表“在惯性系 $\underline{\mathcal{F}}_1$ 中看的导数”、$(\cdot)^{\circ}$ 表“在随体系 $\underline{\mathcal{F}}_2$ 中看的导数”，对任意向量 $\underline{r}$ 有核心的\emph{哥氏关系}
\begin{equation}
  \underline{r}^{\bullet}=\underline{r}^{\circ}+\boldsymbol{\omega}_{21}\times\underline{r}.
  \label{eq:rb-coriolis}
\end{equation}
对其再用一次 \cref{eq:rb-coriolis}（令 $\underline{v}=\underline{r}^{\bullet}$）并合并同类项，得惯性加速度的四项分解：
\begin{equation}
  \underline{r}^{\bullet\bullet}=\underbrace{\underline{r}^{\circ\circ}}_{\text{相对加速度}}
  +\underbrace{2\,\boldsymbol{\omega}_{21}\times\underline{r}^{\circ}}_{\text{哥氏加速度}}
  +\underbrace{\boldsymbol{\omega}_{21}^{\circ}\times\underline{r}}_{\text{角加速度项}}
  +\underbrace{\boldsymbol{\omega}_{21}\times(\boldsymbol{\omega}_{21}\times\underline{r})}_{\text{向心加速度}}.
  \label{eq:rb-coriolis-acc}
\end{equation}
其坐标形式（用 $\mathbf{C}_{12}$ 与机体系角速度 $\boldsymbol{\omega}_2^{21}$）为 $\ddot{\mathbf{r}}_1=\mathbf{C}_{12}[\ddot{\mathbf{r}}_2+2\boldsymbol{\omega}_2^{21\wedge}\dot{\mathbf{r}}_2+\dot{\boldsymbol{\omega}}_2^{21\wedge}\mathbf{r}_2+\boldsymbol{\omega}_2^{21\wedge}\boldsymbol{\omega}_2^{21\wedge}\mathbf{r}_2]$\cite{barfoot2024state}。这正是 IMU 加速度计模型中“杆臂效应”各项的来源。

\paragraph{位姿运动学与广义速度。}位姿的运动学相应为
\begin{equation}
  \dot{\mathbf{T}}=\begin{bmatrix}\boldsymbol{\omega}^\wedge&\boldsymbol{\nu}\\\mathbf{0}^\top&0\end{bmatrix}\mathbf{T}=\boldsymbol{\varpi}^\wedge\mathbf{T},\qquad
  \boldsymbol{\varpi}=\begin{bmatrix}\boldsymbol{\nu}\\\boldsymbol{\omega}\end{bmatrix},
  \label{eq:rb-se3-kin}
\end{equation}
其中广义 $6$-DOF 速度 $\boldsymbol{\varpi}=[\boldsymbol{\nu};\boldsymbol{\omega}]$（\textbf{平移 $\boldsymbol{\nu}$ 在前、转动 $\boldsymbol{\omega}$ 在后}，与本书 $\se(3)$ 排序 $\boldsymbol{\xi}=[\boldsymbol{\rho};\boldsymbol{\phi}]$ 一致）。这个 $\boldsymbol{\varpi}^\wedge$ 正是 \cref{ch:lie} 的 $\se(3)$ 元素\cite{barfoot2024state}。

\paragraph{Frenet--Serret 与独轮车模型（Barfoot 的两个实例）。}给点 $V$ 沿空间曲线运动附一个标架（$1$ 轴切向 $\underline{t}$、$2$ 轴法向 $\underline{n}$、$3$ 轴副法向 $\underline{b}$），随弧长 $s$ 的变化由 \emph{Frenet--Serret 方程}描述（$\kappa$ 曲率、$\tau$ 挠率）：
\begin{equation}
  \frac{d}{ds}\begin{bmatrix}\underline{t}\\\underline{n}\\\underline{b}\end{bmatrix}
  =\begin{bmatrix}0&\kappa&0\\-\kappa&0&\tau\\0&-\tau&0\end{bmatrix}\begin{bmatrix}\underline{t}\\\underline{n}\\\underline{b}\end{bmatrix}.
\end{equation}
乘速度 $v=ds/dt$ 即恢复 Poisson 方程，对应机体角速度 $\boldsymbol{\omega}_v^{vi}=[v\tau,0,v\kappa]^\top$（只有 $2$ 自由度），平移速度 $\boldsymbol{\nu}_v^{vi}=[v,0,0]^\top$。把它代入 \cref{eq:rb-se3-kin}，约束到 $xy$ 平面（令 $\tau=0$、$z$ 分量冻结），得移动机器人最常用的\textbf{独轮车 / 差速驱动模型}：
\begin{equation}
  \dot{x}=v\cos\theta,\qquad \dot{y}=v\sin\theta,\qquad \dot{\theta}=\omega\ (=v\kappa),
  \label{eq:rb-unicycle}
\end{equation}
输入为纵向速度 $v$ 与转速 $\omega$。轮子的\emph{非完整约束}（nonholonomic）使车不能横向平移，只能前滚加转——这是规划与控制章反复出现的模型\cite{barfoot2024state}。

\paragraph{欧拉角速率映射与万向锁的数学面目。}机体角速度 $\boldsymbol{\omega}$ 与欧拉角速率 $\dot{\boldsymbol{\theta}}$ 经一个矩阵 $\mathbf{S}(\boldsymbol{\theta})$ 相联，$\boldsymbol{\omega}=\mathbf{S}(\boldsymbol{\theta})\dot{\boldsymbol{\theta}}$。对 $1$-$2$-$3$（RPY）序列，Barfoot 由 $\dot{\mathbf{C}}_i\mathbf{C}_i^\top=-\mathbf{1}_i^\wedge\dot{\theta}_i$ 与恒等式 (I5) 链式求导得到\cite{barfoot2024state}
\begin{equation}
  \mathbf{S}(\theta_2,\theta_3)=\begin{bmatrix}\cos\theta_2\cos\theta_3&\sin\theta_3&0\\-\cos\theta_2\sin\theta_3&\cos\theta_3&0\\\sin\theta_2&0&1\end{bmatrix},
\end{equation}
反求欧拉角速率要用 $\mathbf{S}^{-1}$，其首行含 $\sec\theta_2$：
\begin{equation}
  \dot{\boldsymbol{\theta}}=\mathbf{S}^{-1}(\boldsymbol{\theta})\,\boldsymbol{\omega}
  =\begin{bmatrix}\sec\theta_2\cos\theta_3&-\sec\theta_2\sin\theta_3&0\\\sin\theta_3&\cos\theta_3&0\\-\tan\theta_2\cos\theta_3&\tan\theta_2\sin\theta_3&1\end{bmatrix}\boldsymbol{\omega},
  \qquad \mathbf{S}^{-1}\xrightarrow{\theta_2\to\pm90^\circ}\infty.
  \label{eq:rb-eulerrate}
\end{equation}
这就是\textbf{万向锁}在运动学上的表现：俯仰接近 $\pm90^\circ$ 时，欧拉角速率趋于无穷，数值崩溃。无奇异的四元数运动学 $\dot{\mathbf{q}}=\tfrac12\mathbf{q}\,[0,\boldsymbol{\omega}]^\top$（或 $\dot{\mathbf{q}}=\tfrac12[0,\boldsymbol{\omega}]^\top\mathbf{q}$，视世界/机体系）则无此问题，故 IMU 积分用四元数而非欧拉角。

\begin{practice}
\begin{enumerate}
  \item （概念）从 \cref{eq:poisson} 出发，说明当 $\boldsymbol{\omega}$ 为常值时 $\mathbf{R}(t)$ 的形式应是什么（提示：这正是 \cref{ch:lie} 的指数映射 $\mathbf{R}(t)=\exp(\boldsymbol{\omega}^\wedge t)\mathbf{R}(0)$）。
  \item （推导）由哥氏关系 \cref{eq:rb-coriolis} 对 $\underline{r}^{\bullet}$ 再求一次导，独立得出四项加速度 \cref{eq:rb-coriolis-acc}，并指出哪一项是“向心”、哪一项是“哥氏”。
  \item 解释为何 $\mathbf{S}^{-1}$ 含 $\sec\theta_2$ 就意味着万向锁，并与 \cref{sec:rb-euler} 的几何描述对应。
  \item （跨章综合，承重）用本章 \cref{eq:poisson} 的 Poisson 方程 + \cref{eq:quat-rot} 的四元数 + \cref{prop:T-inv} 的位姿求逆，写出“给定恒定机体角速度与线速度、把一个机体系点积分到世界系”的完整流程（为 \cref{ch:vio} 的 IMU 预积分铺垫）。
\end{enumerate}
\end{practice}

% ════════════════════════════════════════════════════════════════
\section{（选读）单刚体动力学：SE(3) 上的牛顿--欧拉方程}
\label{sec:rb-dynamics}
\noindent\emph{节首导读：选读（难度 $\bigstar\bigstar\bigstar\bigstar$）。前一节只讲了\emph{运动学}（位姿怎么随速度变）；本节补上\emph{动力学}（速度怎么随受力变），把单刚体的牛顿第二定律统一写成 $\SE(3)$ 上一条紧凑方程，并给出 D'Eleuterio 的广义质量阵。它是含动力学的运动模型（如恒力/恒力矩先验、物体级状态估计）的连续时间过程模型，与 \cref{ch:lie} 的状态转移矩阵配套使用。}

\paragraph{动机：运动学之外还要受力。}\cref{eq:rb-se3-kin} 的位姿运动学 $\dot{\mathbf{T}}=\boldsymbol{\varpi}^\wedge\mathbf{T}$ 只回答“给定广义速度 $\boldsymbol{\varpi}$，位姿怎么变”，把 $\boldsymbol{\varpi}$ 当成已知输入。但真实刚体的 $\boldsymbol{\varpi}$ 由\emph{受力}决定（牛顿第二定律 + 欧拉方程）。把平动与转动的动力学一并写到机体系，恰好能像运动学那样\emph{统一}成 $\SE(3)$ 上的一条向量方程\cite{barfoot2024state}（D'Eleuterio, 1985）。

\begin{definition}[广义质量阵与广义动量]\label{def:rb-genmass}
设刚体质量 $m$、质心在机体系下的坐标 $\mathbf{c}\in\mathbb{R}^3$、关于机体系原点的转动惯量 $\mathbf{I}\in\mathbb{R}^{3\times3}$。定义 $6\times6$ 的\textbf{广义质量阵}（generalized/spatial mass matrix）
\begin{equation}
  \boldsymbol{\mathcal{M}}=\begin{bmatrix}m\mathbf{1}&-m\mathbf{c}^\wedge\\ m\mathbf{c}^\wedge&\mathbf{I}\end{bmatrix}\in\mathbb{R}^{6\times6}.
  \label{eq:rb-genmass}
\end{equation}
它把广义速度 $\boldsymbol{\varpi}=[\boldsymbol{\nu};\boldsymbol{\omega}]$（机体系，平移在前）映为\textbf{广义动量} $\mathbf{p}=\boldsymbol{\mathcal{M}}\boldsymbol{\varpi}=[m(\boldsymbol{\nu}-\mathbf{c}^\wedge\boldsymbol{\omega});\,m\mathbf{c}^\wedge\boldsymbol{\nu}+\mathbf{I}\boldsymbol{\omega}]$，上块是线动量、下块是关于原点的角动量；刚体动能恰为 $T_{\text{kin}}=\tfrac12\boldsymbol{\varpi}^\top\boldsymbol{\mathcal{M}}\boldsymbol{\varpi}$。
\end{definition}

\noindent 非对角块 $\pm m\mathbf{c}^\wedge$ 编码了\emph{质心不在机体系原点}带来的平/转动耦合：若取机体系原点\emph{就在质心}（$\mathbf{c}=\mathbf{0}$），$\boldsymbol{\mathcal{M}}=\operatorname{diag}(m\mathbf{1},\mathbf{I})$ 解耦，退回中学熟悉的“线动量 $m\boldsymbol{\nu}$、角动量 $\mathbf{I}\boldsymbol{\omega}$”。

\begin{figure}[ht]
\centering
\begin{tikzpicture}[>=Stealth,scale=1.0]
  % rigid body blob
  \draw[thick, fill=main!8, draw=main!60!black] plot[smooth cycle] coordinates {(0.2,0.1)(1.4,-0.2)(2.3,0.5)(1.9,1.5)(0.7,1.6)(-0.2,0.9)};
  % body frame at origin O
  \coordinate (O) at (0.4,0.5);
  \fill (O) circle (1.2pt) node[below left] {$O$ (机体原点)};
  \draw[->, thick] (O) -- ++(1.0,0) node[right] {$\underline{1}^x_b$};
  \draw[->, thick] (O) -- ++(0,0.95) node[above] {$\underline{1}^y_b$};
  % center of mass
  \coordinate (C) at (1.4,0.95);
  \fill[second] (C) circle (1.6pt) node[above right] {$\bullet\,\mathbf{c}$ (质心)};
  \draw[->, second, thick] (O) -- (C) node[midway,below right] {$\mathbf{c}$};
  % generalized velocity
  \draw[->, third, very thick] (O) -- ++(-0.9,-0.55) node[below] {$\boldsymbol{\nu}$};
  \draw[->, third, thick] (0.05,1.25) arc (160:30:0.42);
  \node[third] at (0.95,1.75) {$\boldsymbol{\omega}$};
\end{tikzpicture}
\caption{单刚体的机体系 $\{O;\underline{1}^x_b,\underline{1}^y_b\}$、质心偏置 $\mathbf{c}$ 与广义速度 $\boldsymbol{\varpi}=[\boldsymbol{\nu};\boldsymbol{\omega}]$。广义质量阵 \cref{eq:rb-genmass} 的非对角块 $\pm m\mathbf{c}^\wedge$ 正源于 $\mathbf{c}\neq\mathbf{0}$ 的平动--转动耦合；取 $O$ 在质心则解耦。}
\label{fig:rb-genmass}
\end{figure}

\begin{theorem}[$\SE(3)$ 上的单刚体牛顿--欧拉方程]\label{thm:rb-dynamics}
机体系下，刚体的运动学与动力学合写为
\begin{align}
  \text{运动学：}&\quad \dot{\mathbf{T}}=\boldsymbol{\varpi}^\wedge\mathbf{T}, \label{eq:rb-dyn-kin}\\
  \text{动力学：}&\quad \dot{\boldsymbol{\varpi}}=-\boldsymbol{\mathcal{M}}^{-1}\boldsymbol{\varpi}^{\curlywedge\top}\boldsymbol{\mathcal{M}}\boldsymbol{\varpi}+\mathbf{a}, \label{eq:rb-dyn-dyn}
\end{align}
其中 $\mathbf{a}\in\mathbb{R}^6$ 是机体系下的广义外力（每单位质量），$\boldsymbol{\varpi}^\curlywedge$ 是\emph{广义速度的} $6\times6$ 反对称（小伴随）算子
\begin{equation}
  \boldsymbol{\varpi}^\curlywedge=\begin{bmatrix}\boldsymbol{\nu}\\\boldsymbol{\omega}\end{bmatrix}^{\!\curlywedge}=\begin{bmatrix}\boldsymbol{\omega}^\wedge&\boldsymbol{\nu}^\wedge\\\mathbf{0}&\boldsymbol{\omega}^\wedge\end{bmatrix}\in\mathbb{R}^{6\times6},
  \label{eq:rb-curlywedge}
\end{equation}
（此算子即 \cref{ch:lie} 系统引入的小伴随 $\mathrm{ad}$，见 \cref{sec:adjoint}）。
\end{theorem}

\noindent\textbf{逐项读懂动力学方程 \cref{eq:rb-dyn-dyn}。}它就是把欧拉方程“$\mathbf{I}\dot{\boldsymbol{\omega}}+\boldsymbol{\omega}\times\mathbf{I}\boldsymbol{\omega}=\boldsymbol{\tau}$”升级到 $\SE(3)$：$\boldsymbol{\mathcal{M}}\boldsymbol{\varpi}$ 是广义动量，$-\boldsymbol{\varpi}^{\curlywedge\top}\boldsymbol{\mathcal{M}}\boldsymbol{\varpi}$ 是把“动量被机体系转动带着走”产生的\emph{陀螺/离心}项（对应欧拉方程的 $\boldsymbol{\omega}\times\mathbf{I}\boldsymbol{\omega}$ 与平动的 $\boldsymbol{\omega}\times m\boldsymbol{\nu}$），$\mathbf{a}$ 是外力激励。无外力（$\mathbf{a}=\mathbf{0}$）时它描述自由刚体的章动/进动；这套写法的妙处是平动与转动\emph{共用}一个 $\boldsymbol{\varpi}^\curlywedge$，与运动学 \cref{eq:rb-dyn-kin} 的 $\boldsymbol{\varpi}^\wedge$ 形成 $4\times4$/$6\times6$ 的漂亮对仗。

\paragraph{线性化：当 $\delta\mathbf{a}$ 是过程噪声。}状态估计常把广义外力的不确定 $\delta\mathbf{a}$ 当作过程噪声，问它如何经动力学+运动学传到位姿与速度上。围绕标称解扰动 $\mathbf{T}'=\exp(\delta\boldsymbol{\xi}^\wedge)\mathbf{T}$、$\boldsymbol{\varpi}'=\boldsymbol{\varpi}+\delta\boldsymbol{\varpi}$、$\mathbf{a}'=\mathbf{a}+\delta\mathbf{a}$，丢二阶小量后分离出（标称略）线性扰动模型
\begin{equation}
  \begin{bmatrix}\delta\dot{\boldsymbol{\xi}}\\\delta\dot{\boldsymbol{\varpi}}\end{bmatrix}
  =\begin{bmatrix}\boldsymbol{\varpi}^\curlywedge & \mathbf{1}\\ \mathbf{0} & \boldsymbol{\mathcal{M}}^{-1}\big((\boldsymbol{\mathcal{M}}\boldsymbol{\varpi})^{\curlywedge\top}-\boldsymbol{\varpi}^\curlywedge\boldsymbol{\mathcal{M}}\big)\end{bmatrix}
  \begin{bmatrix}\delta\boldsymbol{\xi}\\\delta\boldsymbol{\varpi}\end{bmatrix}
  +\begin{bmatrix}\mathbf{0}\\\delta\mathbf{a}\end{bmatrix},
  \label{eq:rb-dyn-pert}
\end{equation}
其左上块 $\boldsymbol{\varpi}^\curlywedge$ 正是位姿扰动运动学 $\delta\dot{\boldsymbol{\xi}}=\boldsymbol{\varpi}^\curlywedge\delta\boldsymbol{\xi}+\delta\boldsymbol{\varpi}$（与纯运动学线性化一致），右下块来自动力学项对 $\boldsymbol{\varpi}$ 的求导\cite{barfoot2024state}。\cref{eq:rb-dyn-pert} 是一个线性时变随机微分方程（LTV SDE）；其 $12\times12$ 状态转移矩阵一般无闭式、须数值积分，但其结构与 \cref{ch:lie} 的状态转移理论（旋转/位姿分量的解析转移矩阵 $\mathbf{R}(t)\mathbf{R}(s)^\top$、$\boldsymbol{\mathcal{T}}(t)\boldsymbol{\mathcal{T}}(s)^{-1}$）完全衔接，可作为协方差传播的积木。

\begin{insight}[为何把动力学也塞进 $\SE(3)$]
把“位置 + 朝向 + 线速度 + 角速度 + 受力”全部用 $\SE(3)$ 的语言（$\mathbf{T},\boldsymbol{\varpi},\boldsymbol{\mathcal{M}},\boldsymbol{\varpi}^\curlywedge$）写出，好处有三：(1) \emph{无奇异}——姿态始终是 $\mathbf{T}\in\SE(3)$，不会像欧拉角那样万向锁；(2) \emph{平/转统一}——一条 $\dot{\boldsymbol{\varpi}}$ 同时管牛顿与欧拉，省去分别推导；(3) \emph{直接对接估计}——线性化即得 \cref{eq:rb-dyn-pert} 这一标准过程模型，配 \cref{ch:lie} 的状态转移矩阵就能做连续时间的滤波/平滑（如物体级 SLAM、动力学先验的轨迹估计）。这与 \cref{eq:rb-coriolis} 的哥氏分解一脉相承：都是“把惯性系里复杂的受力运动，搬到随体机体系里写简洁”。
\end{insight}

\begin{practice}
\begin{enumerate}
  \item 由 $T_{\text{kin}}=\tfrac12\boldsymbol{\varpi}^\top\boldsymbol{\mathcal{M}}\boldsymbol{\varpi}$ 与 \cref{eq:rb-genmass} 展开验证：当 $\mathbf{c}=\mathbf{0}$ 时退化为 $\tfrac12 m\|\boldsymbol{\nu}\|^2+\tfrac12\boldsymbol{\omega}^\top\mathbf{I}\boldsymbol{\omega}$。
  \item （概念）说明 \cref{eq:rb-dyn-dyn} 的陀螺项 $-\boldsymbol{\varpi}^{\curlywedge\top}\boldsymbol{\mathcal{M}}\boldsymbol{\varpi}$ 在 $\mathbf{c}=\mathbf{0}$、纯转动（$\boldsymbol{\nu}=\mathbf{0}$）时如何还原为经典欧拉方程的 $-\boldsymbol{\omega}\times\mathbf{I}\boldsymbol{\omega}$。\emph{答案线索}：用 \cref{eq:rb-curlywedge} 的下块 $\boldsymbol{\omega}^\wedge$ 与 $(\boldsymbol{\omega}^\wedge)^\top=-\boldsymbol{\omega}^\wedge$。
  \item （跨章）指出 \cref{eq:rb-dyn-pert} 的左上块 $\boldsymbol{\varpi}^\curlywedge$ 为何恰是位姿扰动运动学的系统矩阵（对照 \cref{ch:lie} 的位姿扰动方程）。
\end{enumerate}
\end{practice}

% ════════════════════════════════════════════════════════════════
\section{（选读）相似、仿射与射影变换}
\label{sec:rb-projective}
\noindent\emph{节首导读：选读（服务相机模型章）。欧氏变换保长度夹角；逐级放松约束得到更一般的变换层级。}

欧氏变换最简单，保长度、夹角、体积；放松约束得到更一般的变换层级，相机成像章会用到\cite{gaoxiang2019slam14}。
\begin{align}
  \text{相似}\ \mathbf{T}_S&=\begin{bmatrix}s\mathbf{R}&\mathbf{t}\\\mathbf{0}^\top&1\end{bmatrix}\quad(7\text{ 自由度，允许均匀缩放 }s),\\
  \text{仿射}\ \mathbf{T}_A&=\begin{bmatrix}\mathbf{A}&\mathbf{t}\\\mathbf{0}^\top&1\end{bmatrix}\quad(\mathbf{A}\text{ 可逆，}12\text{ 自由度，保平行}),\\
  \text{射影}\ \mathbf{T}_P&=\begin{bmatrix}\mathbf{A}&\mathbf{t}\\\mathbf{a}^\top&v\end{bmatrix}\quad(15\text{ 自由度，最一般}).
\end{align}
不变性质逐级减弱，且\emph{从上到下有包含关系}：
\begin{table}[H]\centering\small
\caption{常见变换的性质比较（十四讲表 3-1）}
\label{tab:rb-transforms}
\begin{tabular}{lccl}
\toprule
变换 & 矩阵形式 & 自由度 & 不变性质 \\
\midrule
欧氏 & $[\mathbf{R},\mathbf{t};\mathbf{0}^\top,1]$ & $6$ & 长度、夹角、体积 \\
相似 & $[s\mathbf{R},\mathbf{t};\mathbf{0}^\top,1]$ & $7$ & 体积比 \\
仿射 & $[\mathbf{A},\mathbf{t};\mathbf{0}^\top,1]$ & $12$ & 平行性、体积比 \\
射影 & $[\mathbf{A},\mathbf{t};\mathbf{a}^\top,v]$ & $15$ & 接触平面的相交与相切（共线性、交比）\\
\bottomrule
\end{tabular}
\end{table}

\noindent\textbf{真实世界到相机照片的变换是射影变换}（近大远小，方形地砖在像里成不规则四边形）；相机焦距趋于无穷时退化为仿射——这条线索留给 \cref{ch:camera}。相似变换群 $\mathrm{Sim}(3)$ 在单目 SLAM 的尺度处理中重要（其李代数见 \cref{ch:lie} 的 $\mathrm{Sim}(3)$ 一节）。Barfoot 在同一框架下顺势导出了相机内参矩阵 $\mathbf{K}$、本质矩阵 $\mathbf{E}=\mathbf{C}_{ba}^\top\mathbf{r}_b^{ab\wedge}$、基础矩阵 $\mathbf{F}=\mathbf{K}_a^{-\top}\mathbf{E}\mathbf{K}_b^{-1}$ 与单应矩阵 $\mathbf{H}_{ba}=\tfrac{z_a}{z_b}\mathbf{C}_{ba}(\mathbf{I}+\tfrac{1}{d_a}\mathbf{r}_a^{ba}\mathbf{n}_a^\top)$ 等射影几何量\cite{barfoot2024state}，它们的系统展开在 \cref{ch:camera} 与对极几何相关章节进行。

% ════════════════════════════════════════════════════════════════
\section{实践：Eigen 几何模块}
\label{sec:rb-eigen}
\noindent\emph{节首导读：骨干配套。把前面的数学落到可运行代码，并完整复现一道数值例。}

\paragraph{桥接：从纸面到代码。}前面十节我们在纸上把旋转、位姿、各种表示及其互转推了个遍。但教材若止于公式，读者很难真正“相信”这些抽象对象——它们到底长什么样、算起来对不对？本节把全部理论落到一个能编译运行的库 \textbf{Eigen} 上：先用它最基础的\emph{矩阵模块}熟悉数据类型与运算（呼应 \cref{sec:rb-vector}--\cref{sec:rb-se3} 的线性代数底座），再用它的\emph{几何模块}演示旋转/位姿的表示与互转（呼应 \cref{sec:rb-axisangle}--\cref{sec:rb-quat}），最后完整复现一道坐标变换数值例。本节分\emph{两段}实践——矩阵模块与几何模块——缺一不可\cite{gaoxiang2019slam14}。

Eigen 是 C++ 头文件线性代数库（g2o、Sophus、Ceres 都依赖它），其 Geometry 模块直接提供各旋转/位姿表示及互转。Eigen 是\emph{纯头文件库}（无 \texttt{.so}/\allowbreak\texttt{.a}），使用时只需 \texttt{include}，无需链接（CMake 里只给 \texttt{include\_directories} 而\emph{不}写 \texttt{target\_link\_libraries}）\cite{gaoxiang2019slam14}。

\paragraph{第一段实践：Eigen 矩阵模块（十四讲 \texttt{eigenMatrix.\allowbreak cpp}，全录可复现）。}在用几何模块表示旋转之前，先打好地基：Eigen 里\emph{一切}向量与矩阵都是模板类 \texttt{Eigen::Matrix<类型, 行, 列>}，三维向量、$3\times3$ 旋转矩阵都是它的特例。下面这段代码演示矩阵的声明、初始化、逐元素访问、转置/迹/逆/行列式、特征值分解，以及用矩阵分解高效求解线性方程 $\mathbf{A}\mathbf{x}=\mathbf{b}$——这些正是后续所有数值实现的“砖块”。

\begin{codebox}[C++]{Eigen 矩阵模块:类型、基本运算与解方程}
#include <iostream>
#include <ctime>
#include <Eigen/Core>     // Eigen 核心（矩阵/向量）
#include <Eigen/Dense>    // 稠密矩阵代数（逆、特征值、分解等）
using namespace Eigen;

#define MATRIX_SIZE 50

int main() {
    // Matrix<类型, 行, 列>:所有向量/矩阵的统一模板
    Matrix<float, 2, 3> matrix_23;          // 2x3 float 矩阵
    Vector3d v_3d;                          // = Matrix<double,3,1>，三维列向量
    Matrix<float, 3, 1> vd_3d;              // 与 Vector3f 等价
    Matrix3d matrix_33 = Matrix3d::Zero();  // 3x3，初始化为零
    MatrixXd matrix_x;                      // 动态大小（编译期未知）

    matrix_23 << 1, 2, 3, 4, 5, 6;          // 逗号初始化（行优先填入）
    for (int i = 0; i < 2; i++)             // 用 () 逐元素访问
        for (int j = 0; j < 3; j++) std::cout << matrix_23(i, j) << "\t";

    v_3d << 3, 2, 1;
    vd_3d << 4, 5, 6;
    // Eigen 不做隐式类型提升:float 矩阵乘 double 向量必须显式 cast
    Matrix<double, 2, 1> result = matrix_23.cast<double>() * v_3d;
    Matrix<float, 2, 1> result2 = matrix_23 * vd_3d;   // 同类型则可直接乘

    matrix_33 = Matrix3d::Random();         // 随机矩阵
    std::cout << matrix_33.transpose();     // 转置
    std::cout << matrix_33.sum();           // 元素和
    std::cout << matrix_33.trace();         // 迹（呼应轴角求转角用到的 tr R）
    std::cout << 10 * matrix_33;            // 数乘
    std::cout << matrix_33.inverse();       // 逆
    std::cout << matrix_33.determinant();   // 行列式（SO(3) 判定要 det = 1）

    // 实对称矩阵的特征值分解（保证可对角化）
    SelfAdjointEigenSolver<Matrix3d> es(matrix_33.transpose() * matrix_33);
    std::cout << es.eigenvalues() << es.eigenvectors();

    // 解方程 A x = b:直接求逆量大，矩阵分解更快
    Matrix<double, MATRIX_SIZE, MATRIX_SIZE> A
        = MatrixXd::Random(MATRIX_SIZE, MATRIX_SIZE);
    A = A * A.transpose();                   // 构造半正定（对称）矩阵
    Matrix<double, MATRIX_SIZE, 1> b = MatrixXd::Random(MATRIX_SIZE, 1);

    clock_t t = clock();
    Matrix<double, MATRIX_SIZE, 1> x = A.inverse() * b;          // 1) 直接求逆
    std::cout << 1000.0 * (clock() - t) / CLOCKS_PER_SEC << "ms";
    t = clock();
    x = A.colPivHouseholderQr().solve(b);                        // 2) QR 分解
    std::cout << 1000.0 * (clock() - t) / CLOCKS_PER_SEC << "ms";
    t = clock();
    x = A.ldlt().solve(b);                                       // 3) Cholesky(LDLT)
    std::cout << 1000.0 * (clock() - t) / CLOCKS_PER_SEC << "ms";
    return 0;
}
\end{codebox}
\noindent\textbf{读这段代码的要点}（十四讲实践注记的精炼）：(1) \emph{固定大小} vs \emph{动态大小}——编译期已知尺寸的 \texttt{Matrix3d}、\texttt{Vector3d} 比 \texttt{MatrixXd} 快，旋转/变换矩阵应一律用固定大小。(2) \emph{无自动类型提升}——\texttt{float} 与 \texttt{double} 矩阵相乘必须显式 \texttt{cast<double>()}，否则 Eigen 抛出大写错误 “YOU MIXED DIFFERENT NUMERIC TYPES”；维度不匹配则报 “YOU MIXED MATRICES OF DIFFERENT SIZES”——出错时直接搜大写部分定位。(3) 三种解方程方式在数值上等价，但\emph{效率差异显著}：直接求逆 $O(n^3)$ 且数值稳定性差，QR、Cholesky 分解通常快得多（对正定矩阵 \texttt{ldlt} 最快）；这条经验在 \cref{ch:nlopt}、\cref{ch:slam_est} 求解大规模正规方程时反复用到。(4) \texttt{trace} 与 \texttt{determinant} 不只是 API——它们正是本章判定 $\SO(3)$（$\det=1$）与从矩阵反求转角（\cref{eq:rb-trace} 的 $\operatorname{tr}\mathbf{R}$）所需的运算。

\paragraph{第二段实践：Eigen 几何模块。}地基打好，现在用几何模块直接操作旋转与位姿——它把轴角、旋转矩阵、四元数、欧氏变换都封装成专门类型，并重载了运算符，使“旋转一个向量”“复合两个位姿”读起来就像数学公式。

\begin{codebox}[C++]{Eigen:各旋转表示的构造与互转、坐标变换}
#include <Eigen/Core>
#include <Eigen/Geometry>
using namespace Eigen;

// 旋转:轴角 -> 旋转矩阵 / 四元数（沿 Z 轴 45 度）
AngleAxisd rv(M_PI / 4, Vector3d(0, 0, 1));
Matrix3d   R = rv.toRotationMatrix();         // 轴角 -> 矩阵
Quaterniond q(rv);                            // 轴角 -> 四元数（Hamilton）
q = Quaterniond(R);                           // 矩阵 -> 四元数
Vector3d euler = R.eulerAngles(2, 1, 0);      // 矩阵 -> 欧拉角(ZYX = yaw,pitch,roll)

// 用各表示旋转同一个向量，结果一致:
Vector3d v(1, 0, 0);
Vector3d v1 = R * v;                           // 矩阵作用
Vector3d v2 = q * v;                           // 四元数作用（内部即 q v q^-1）

// 位姿 T_wc（相机到世界）:Isometry3d 即 4x4 欧氏变换
Isometry3d Twc = Isometry3d::Identity();
Twc.rotate(q);                                 // 设旋转 R_wc
Twc.pretranslate(Vector3d(1, 3, 4));           // 设平移 t_wc（相机在世界中的位置）
Vector3d p_w = Twc * Vector3d(0, 0, 0);        // 相机系原点 -> 世界 = t_wc
Isometry3d Tcw = Twc.inverse();                // 世界到相机（整体求逆，勿对 t 取负！）

// 帧链:已知 T_wc1, T_wc2，求点从 c1 系到 c2 系
// p_c2 = T_c2,w * T_w,c1 * p_c1 = (Twc2.inverse()) * Twc1 * p_c1
\end{codebox}

\paragraph{Eigen 中各表示对应的数据类型（速查，写代码常回看）。}下表把“数学对象 $\to$ Eigen 类型”一一对照，写代码时极常翻查（以双精度 \texttt{d} 结尾，改 \texttt{f} 即单精度，二者同样\emph{不能}自动互转）\cite{gaoxiang2019slam14}：
\begin{table}[H]\centering\small
\caption{旋转/位姿各表示在 Eigen 中对应的数据类型}
\label{tab:rb-eigen-types}
\begin{tabular}{lll}
\toprule
数学对象 & Eigen 类型 & 维度 / 备注 \\
\midrule
旋转矩阵 & \texttt{Eigen::Matrix3d} & $3\times3$ \\
旋转向量（轴角） & \texttt{Eigen::AngleAxisd} & 底层非矩阵，重载运算符可当矩阵用 \\
欧拉角 & \texttt{Eigen::Vector3d} & $3\times1$，由 \texttt{eulerAngles(2,1,0)} 得 ZYX \\
四元数 & \texttt{Eigen::Quaterniond} & 构造按 $(w,x,y,z)$；\texttt{coeffs()} 存 $(x,y,z,w)$ \\
欧氏变换（位姿） & \texttt{Eigen::Isometry3d} & 名为 3D，实为 $4\times4$ \\
仿射变换 & \texttt{Eigen::Affine3d} & $4\times4$ \\
射影变换 & \texttt{Eigen::Projective3d} & $4\times4$ \\
\bottomrule
\end{tabular}
\end{table}
\noindent\textbf{务必记住的两处“坑”}：四元数构造参数顺序 $(w,x,y,z)$ 与内存 \texttt{coeffs()} 顺序 $(x,y,z,w)$ \emph{不同}（呼应 \cref{def:quat} 后的库差异说明）；用同一旋转分别经矩阵、轴角、四元数三条路去旋转同一向量 $\mathbf{v}$，结果\emph{必须一致}（不一致就说明哪里写错了），这也是验证自己理解是否到位的最快自测。

\paragraph{完整数值例（十四讲“小萝卜”坐标变换，全录可复现）。}世界系 $W$，机器人一号、二号坐标系 $R_1,R_2$。一号位姿 $\mathbf{q}_1=[0.35,0.2,0.3,0.1]$、$\mathbf{t}_1=[0.3,0.1,0.1]^\top$；二号位姿 $\mathbf{q}_2=[-0.5,0.4,-0.1,0.2]$、$\mathbf{t}_2=[-0.1,0.5,0.3]^\top$。这里 $\mathbf{q},\mathbf{t}$ 表达的是\emph{世界系到机器人系}的变换 $\mathbf{T}_{R_kW}$。一号看到某点在自身系坐标 $\mathbf{p}_{R_1}=[0.5,0,0.2]^\top$，求该点在二号系下的坐标。按相邻相消，$\mathbf{p}_{R_2}=\mathbf{T}_{R_2W}\mathbf{T}_{WR_1}\mathbf{p}_{R_1}=\mathbf{T}_{R_2W}\mathbf{T}_{R_1W}^{-1}\mathbf{p}_{R_1}$。

\begin{codebox}[C++]{坐标变换数值例:求 p 在二号机器人系下坐标}
#include <Eigen/Core>
#include <Eigen/Geometry>
using namespace Eigen;

int main() {
    // 构造参数顺序为 (w,x,y,z);务必归一化
    Quaterniond q1(0.35, 0.2, 0.3, 0.1), q2(-0.5, 0.4, -0.1, 0.2);
    q1.normalize();
    q2.normalize();
    Vector3d t1(0.3, 0.1, 0.1), t2(-0.1, 0.5, 0.3);
    Vector3d p1(0.5, 0, 0.2);

    Isometry3d T1w(q1), T2w(q2);   // T_{R1,W}, T_{R2,W}
    T1w.pretranslate(t1);
    T2w.pretranslate(t2);

    Vector3d p2 = T2w * T1w.inverse() * p1;   // p_{R2} = T2w * (T1w)^-1 * p1
    // 输出: -0.0309731  0.73499  0.296108
    return 0;
}
\end{codebox}
\noindent 运行结果 $\mathbf{p}_{R_2}=[-0.0309731,\ 0.73499,\ 0.296108]^\top$。\textbf{关键}：四元数使用前务必 \texttt{normalize()}；$\mathbf{T}_{WR_1}=\mathbf{T}_{R_1W}^{-1}$ 用整体 \texttt{inverse()} 而非分别取负（呼应 \cref{prop:T-inv} 与平移取负陷阱）。

\paragraph{轨迹可视化的帧含义。}轨迹文件每行常存 \texttt{time, $t_x,t_y,t_z,q_x,q_y,q_z,q_w$}，记录\emph{世界系到机器人系}（或其逆）的位姿；画出的轨迹点严格说是机器人系原点在世界系的坐标 $\mathbf{O}_W=\mathbf{T}_{WR}\mathbf{O}_R=\mathbf{t}_{WR}$（$\mathbf{O}_R=\mathbf{0}$）——即 $\mathbf{T}_{WR}$ 的平移块，这正是存 $\mathbf{T}_{WR}$ 更直观的原因（\cref{sec:rb-convention} 的洞见）。要把轨迹画出来：逐行读入 $(\mathbf{t},\mathbf{q})$ 拼成 \texttt{Isometry3d}，对每个位姿取原点 $\mathbf{O}_W=\mathbf{T}_{WR}\,\mathbf{0}$ 与三个轴端点 $\mathbf{T}_{WR}(\ell\,\mathbf{e}_i)$（$\ell$ 取一小段长度），用红/绿/蓝三色线段画出该位姿的 $X/Y/Z$ 坐标轴，再用一条线把相邻位姿的原点连起来即得轨迹（Linux 下常用基于 OpenGL 的 Pangolin，Python 里用 Matplotlib 的 \texttt{plot3D} 同理）\cite{gaoxiang2019slam14}。三色坐标轴让人一眼看出每处的\emph{朝向}，连线给出\emph{路径}，二者叠加即“带姿态的轨迹”。

\begin{note}
\textbf{实务要点.}\;(1) Eigen 是纯头文件库，CMake 中只需 \texttt{include\_directories}，\emph{无需} \texttt{target\_link\_libraries}。(2) Eigen 矩阵\emph{不支持自动类型提升}：\texttt{float} 与 \texttt{double} 矩阵相乘要显式 \texttt{cast<double>()}，否则报 “YOU MIXED DIFFERENT NUMERIC TYPES”；维度不匹配报 “YOU MIXED MATRICES OF DIFFERENT SIZES”——出错时直接找大写部分定位。(3) 编译期已知大小的矩阵（\texttt{Matrix3d} 等）比动态大小（\texttt{MatrixXd}）快，旋转/变换矩阵应用固定大小。(4) 用 \texttt{std::vector} 存固定大小 Eigen 类型须用 \texttt{Eigen::aligned\_allocator}（内存对齐），漏写会崩溃。(5) \texttt{eulerAngles(2,1,0)} 给 ZYX（yaw-pitch-roll）顺序，其万向锁会让输出跳变，调试用即可。(6) 求解 $\mathbf{A}\mathbf{x}=\mathbf{b}$ 直接求逆运算量大，用 QR（\texttt{colPivHouseholderQr}）或 Cholesky（\texttt{ldlt}）分解更快。
\end{note}

\paragraph{块操作：取子矩阵。}Eigen 用 \texttt{.\allowbreak block<r,c>(i,j)} 取出从 $(i,j)$ 起的 $r\times c$ 子矩阵，它是\emph{可读可写}的视图——既能读子块，也能直接赋值改写原矩阵对应区域。下面把一个大矩阵的左上角 $3\times3$ 块整体置为单位阵 $\mathbf{I}_3$，正是从大状态矩阵里“挖出”一块姿态/协方差子块来操作的常见手法（\cref{ch:nlopt} 装配雅可比、\cref{ch:slam_est} 提取协方差子块都靠它）：
\begin{codebox}[C++]{Eigen 块操作:取左上角 3x3 块并赋值为单位阵}
#include <Eigen/Core>
using namespace Eigen;

MatrixXd M = MatrixXd::Random(6, 6);          // any large matrix
M.block<3, 3>(0, 0) = Matrix3d::Identity();   // set top-left 3x3 block to I_3
// dynamic-size form (block dims given at runtime):
// M.block(0, 0, 3, 3) = Matrix3d::Identity();
// read-only access: Matrix3d topleft = M.block<3, 3>(0, 0);
\end{codebox}

\begin{practice}
\begin{enumerate}
  \item （块操作）取一个大矩阵（如 $6\times6$ 随机阵）的左上角 $3\times3$ 块，赋值为 $\mathbf{I}_3$；再把右上角 $3\times3$ 块读出并打印。\emph{答案线索}：\texttt{M.\allowbreak block<3,3>(0,0)=Matrix3d::Identity();}，可读可写。
  \item （交互体验，强烈建议亲手跑）写一个小程序：维护一个旋转 $\mathbf{R}$（或等价四元数 $\mathbf{q}$），用键盘/鼠标输入一个\emph{增量}旋转 $\Delta\mathbf{R}=\exp(\delta\boldsymbol{\phi}^\wedge)$（如方向键各对应绕某轴的小角增量），每步令 $\mathbf{q}\leftarrow\mathbf{q}\,\Delta\mathbf{q}$ 并\emph{实时打印}该姿态的四种表示——旋转矩阵 \texttt{R}、平移 $\mathbf{t}$、欧拉角 \texttt{R.\allowbreak eulerAngles(2,1,0)}、四元数 \texttt{q.coeffs()}（必要时可用 Pangolin 把坐标轴画出来联动显示）。亲手转几下后体会：\emph{除欧拉角外，其余三种表示的数值对人都不直观}（看着一串数想象不出朝向），但欧拉角在大俯仰角处会跳变——这正印证了 \cref{sec:rb-euler} 的奇异性与 \cref{sec:rb-compare} 的选型结论。\emph{自检}：用同一 $\mathbf{q}$ 经矩阵、轴角、四元数三条路旋转同一向量，结果应一致。
  \item （轨迹绘制）按“轨迹可视化的帧含义”一段的指引，读入若干 $(\mathbf{t},\mathbf{q})$ 位姿，对每个位姿画红/绿/蓝三色坐标轴并把相邻原点连线，画出一条带姿态的三维轨迹（Pangolin 或 Matplotlib 任选）。
\end{enumerate}
\end{practice}

% ════════════════════════════════════════════════════════════════
\section*{本章常见误解汇总}
\begin{table}[H]\centering\small
\begin{tabular}{p{0.46\textwidth} p{0.46\textwidth}}
\toprule
\textbf{常见误解} & \textbf{正确理解} \\
\midrule
向量就是它的坐标 & 向量客观存在；坐标依赖所选基（\cref{eq:coord}）\\
$\mathbf{R}^{-1}=\mathbf{R}^\top$ 是巧合 & 正交性定义后果，旋转保范的必然（\cref{prop:R-inv}）\\
$\mathbf{R}^\top\mathbf{R}=\mathbf{I}$ 就是旋转 & 还需 $\det=+1$，否则可能是镜像（\cref{thm:rb-orthonormal}）\\
$\mathbf{t}_{wc}=-\mathbf{t}_{cw}$ & 实为 $\mathbf{t}_{cw}=-\mathbf{R}_{wc}^\top\mathbf{t}_{wc}$（\cref{eq:T-inv}）\\
Handbook 的 $\mathbf{R}_a^b$ 与本书 $\mathbf{R}_{ab}$ 同向 & 相反！$\mathbf{R}_a^b=\mathbf{R}_{ba}$（\cref{tab:rb-frame-notation}）\\
欧拉角能无奇异表所有旋转 & 万向锁；三参数必有奇异（\cref{sec:rb-compare}）\\
$\mathbf{q}$ 与 $-\mathbf{q}$ 是不同旋转 & 双重覆盖，二者同一旋转 \\
Hamilton 与 JPL 四元数可混用 & $\mathrm{ij}=\pm\mathrm{k}$ 相反，复合顺序相反，必出错 \\
Rodrigues 第三项符号无所谓 & $\pm\sin\theta\,\mathbf{a}^\wedge$ 区分主动/被动、左右手，差一个转置 \\
\bottomrule
\end{tabular}
\end{table}

% ════════════════════════════════════════════════════════════════
\section*{本章小结}

\begin{table}[H]\centering\small
\caption{符号表}
\begin{tabular}{lll}
\toprule
符号 & 含义 & 首次出现 \\
\midrule
$(\cdot)^\wedge,(\cdot)^\vee$ & 向量$\to$反对称阵 / 其逆 & \cref{eq:cross} \\
$\mathbf{R}\in\SO(3)$ & 旋转矩阵（正交，$\det1$） & \cref{def:so3-rep} \\
$\mathbf{R}_{wc},\mathbf{T}_{wc}$ & 相机系$\to$世界系的旋转/位姿 & \cref{sec:rb-convention} \\
$\mathbf{T}\in\SE(3)$ & 变换矩阵（齐次） & \cref{def:se3-rep} \\
$\boldsymbol{\phi}=\theta\mathbf{a}$ & 旋转向量（轴角） & \cref{sec:rb-axisangle} \\
$\mathbf{q}=[w,\mathbf{v}]$ & 四元数（Hamilton，实部在前） & \cref{def:quat} \\
$\mathbf{q}^{+},\mathbf{q}^{\oplus}$ & 四元数左乘/右乘矩阵 & \cref{eq:rb-quat-lr} \\
$\{\boldsymbol{\varepsilon},\eta\}$ & Euler 参数（矢部/标部） & \cref{eq:quat2axis} \\
$\mathbf{g}=\mathbf{a}\tan\tfrac{\theta}{2}$ & Gibbs 向量 & \cref{eq:rb-gibbs} \\
$\boldsymbol{\omega},\boldsymbol{\varpi}$ & 角速度 / 含线速度的广义速度 & \cref{eq:poisson} \\
$\mathbf{S}(\boldsymbol{\theta})$ & 欧拉角速率映射 & \cref{eq:rb-eulerrate} \\
\bottomrule
\end{tabular}
\end{table}

\begin{table}[H]\centering\small
\caption{定理速查表}
\begin{tabular}{lll}
\toprule
公式 & 一句话说明 & 对应 \\
\midrule
反对称恒等式 (I1)--(I6) & 把外积/旋转写成可机械化简的矩阵代数 & \cref{eq:rb-id3} \\
旋转矩阵导出 \cref{eq:R-derive} & 基变换的内积即方向余弦 & \cref{sec:rb-rotation} \\
正交性 \cref{thm:rb-orthonormal} & 旋转矩阵列为标准正交基，$\det=+1$ & \cref{sec:rb-rotation}\\
$\SE(3)$ 逆 \cref{eq:T-inv} & 平移项 $-\mathbf{R}^\top\mathbf{t}$，非 $-\mathbf{t}$ & \cref{prop:T-inv} \\
Rodrigues \cref{eq:rodrigues-rb} & 轴角$\to$矩阵（几何+级数两证） & \cref{thm:rodrigues-rb} \\
迹求转角 \cref{eq:rb-trace} & $\theta=\arccos\frac{\operatorname{tr}\mathbf{R}-1}{2}$ & \cref{sec:rb-axisangle} \\
四元数$\to$矩阵 \cref{eq:quat2R} & $\mathbf{R}=\mathbf{v}\mathbf{v}^\top+s^2\mathbf{I}+2s\mathbf{v}^\wedge+(\mathbf{v}^\wedge)^2$ & \cref{thm:quat2R} \\
四元数$\to$轴角 \cref{eq:quat2axis} & $\theta=2\arccos s$，双重覆盖 & \cref{sec:rb-quat} \\
Gibbs 两式 \cref{eq:rb-cayley} & 无三角函数的旋转矩阵，Rodrigues 代数证 & \cref{sec:rb-gibbs} \\
Poisson \cref{eq:poisson} & $\dot{\mathbf{R}}=\boldsymbol{\omega}^\wedge\mathbf{R}$，伏笔李代数 & \cref{sec:rb-kinematics} \\
独轮车 \cref{eq:rb-unicycle} & 平面差速模型，非完整约束 & \cref{sec:rb-kinematics} \\
欧拉速率 \cref{eq:rb-eulerrate} & $\mathbf{S}^{-1}$ 含 $\sec\theta_2$ = 万向锁 & \cref{sec:rb-kinematics} \\
广义质量阵 \cref{eq:rb-genmass} & $\boldsymbol{\mathcal{M}}=[m\mathbf{1},-m\mathbf{c}^\wedge;m\mathbf{c}^\wedge,\mathbf{I}]$，动能 $\tfrac12\boldsymbol{\varpi}^\top\boldsymbol{\mathcal{M}}\boldsymbol{\varpi}$ & \cref{def:rb-genmass} \\
$\SE(3)$ 牛顿--欧拉 \cref{eq:rb-dyn-dyn} & $\dot{\boldsymbol{\varpi}}=-\boldsymbol{\mathcal{M}}^{-1}\boldsymbol{\varpi}^{\curlywedge\top}\boldsymbol{\mathcal{M}}\boldsymbol{\varpi}+\mathbf{a}$ & \cref{thm:rb-dynamics} \\
\bottomrule
\end{tabular}
\end{table}

\begin{table}[H]\centering\small
\caption{知识点总表}
\begin{tabular}{cllcl}
\toprule
\# & 知识点 & 核心要点 & 难度 & 脉络层 \\
\midrule
1 & 向量与坐标、$(\cdot)^\wedge$ & 向量客观、坐标依赖基；外积$\to$反对称阵；恒等式全套 & \(\bigstar\) & 骨干\\
2 & 旋转矩阵 & 基变换导出；$\SO(3)$ 正交约束\emph{证明} & \(\bigstar\bigstar\) & 骨干\\
3 & 齐次坐标 / $\SE(3)$ & 旋转+平移变线性；逆有闭式 & \(\bigstar\bigstar\) & 骨干\\
4 & 记号约定 & 三书帧记号对照；$\mathbf{t}_{wc}\neq-\mathbf{t}_{cw}$ & \(\bigstar\bigstar\) & 骨干\\
5 & 轴角 / Rodrigues & 几何+级数两向证明；与李代数互印 & \(\bigstar\bigstar\bigstar\) & 骨干\\
6 & 欧拉角 / 万向锁 & 直观但有奇异；三参数必奇异 & \(\bigstar\bigstar\) & 次优\\
7 & 四元数 & Hamilton 全套运算；旋转点；双重覆盖 & \(\bigstar\bigstar\bigstar\) & 骨干\\
8 & Gibbs / 各表示总览 & Cayley 参数；Rodrigues 代数两证 & \(\bigstar\bigstar\bigstar\) & 选读\\
9 & 运动学 & Poisson；哥氏；Frenet--Serret；独轮车；欧拉速率发散 & \(\bigstar\bigstar\bigstar\) & 次优\\
10 & 单刚体动力学 & 广义质量阵 $\boldsymbol{\mathcal{M}}$；$\SE(3)$ 牛顿--欧拉；过程模型线性化 & \(\bigstar\bigstar\bigstar\bigstar\) & 选读\\
11 & 射影变换层级 & 欧氏$\to$相似$\to$仿射$\to$射影 & \(\bigstar\bigstar\) & 选读\\
\bottomrule
\end{tabular}
\end{table}

\section*{延伸阅读}
\begin{itemize}
  \item \textbf{几何直觉与代码}：视觉 SLAM 十四讲\cite{gaoxiang2019slam14} 第 3 章——本章叙事主线、四元数运算与 Eigen/Pangolin 实践来源。
  \item \textbf{严谨 frame 体系}：Barfoot\cite{barfoot2024state} 第 6 章——vectrix 记号、主轴/Euler 参数/Gibbs、哥氏加速度、Poisson、Frenet--Serret、独轮车模型、相机/立体/RAE/IMU 传感器模型（其中 RAE 即距离-方位-俯仰观测模型，本书在 \cref{sec:lidar-rae} 给出其雅可比与退化分析）。
  \item \textbf{现代记号与表示选择}：SLAM Handbook\cite{carlone2026handbook} 第 2 章 / 记号表——$\mathbf{R}_a^b$ 约定、$\mathrm{Exp}/\mathrm{Log}$ 捷径、过参数化与流形动机（其李群优化部分见 \cref{ch:lie}）。
  \item \textbf{“三参数必有奇异”}：Stuelpnagel\cite{stuelpnagel1964parametrization}（旋转群参数化的拓扑障碍）；统一李理论的现代综述见 Solà 等\cite{sola2018micro}。
\end{itemize}

\section*{本章与后续章节的关系}
\begin{table}[H]\centering\small
\begin{tabular}{lll}
\toprule
后续章节 & 关系 & 本章铺垫 \\
\midrule
\cref{ch:lie} 李群与李代数 & 旋转/位姿的\emph{估计与优化} & $\SO(3)/\SE(3)$、Rodrigues、Poisson、广义速度 \\
\cref{ch:camera} 相机模型 & 重投影、射影变换、内参 & 齐次坐标、位姿、$\mathbf{T}_{cw}$、$\mathbf{K}/\mathbf{E}/\mathbf{F}/\mathbf{H}$ \\
\cref{ch:vo} 视觉里程计 & 位姿求解 PnP/ICP & 旋转表示、坐标变换 \\
\cref{ch:vio} VIO & IMU 姿态积分、预积分 & 四元数、Poisson、哥氏加速度 \\
\bottomrule
\end{tabular}
\end{table}

\section*{故障排查手册}
\begin{table}[H]\centering\small
\begin{tabular}{p{0.22\textwidth} p{0.26\textwidth} p{0.34\textwidth} l}
\toprule
症状 & 可能原因 & 排查步骤 & 相关 \\
\midrule
坐标变换结果方向/位置完全错 & 帧下标方向写反，或平移直接取负 & 核对左目标右源 vs Handbook 下源上目标；位姿求逆用 \cref{eq:T-inv} 整体算 & \cref{sec:rb-convention} \\
两姿态四元数插值/求差走“远路” & 双重覆盖，$\mathbf{q}/-\mathbf{q}$ 未统一符号 & 比较前令 $w\ge0$ & \cref{sec:rb-quat} \\
旋转矩阵转出的欧拉角突然跳变 & 俯仰接近 $\pm90^\circ$ 的万向锁 & 内部改用四元数/矩阵，欧拉角仅显示 & \cref{sec:rb-euler} \\
四元数旋转向量结果错或 \texttt{NaN} & 未归一化；Hamilton/JPL 或分量序混淆 & \texttt{normalize}；核对约定与 $(w,x,y,z)$ vs $(x,y,z,w)$ & \cref{sec:rb-quat} \\
\texttt{std::vector<Isometry3d>} 崩溃 & Eigen 固定大小类型内存对齐 & 用 \texttt{Eigen::aligned\_allocator} & \cref{sec:rb-eigen} \\
\bottomrule
\end{tabular}
\end{table}

% ════════════════════════════════════════════════════════════════
\section*{推导附录}
\label{sec:rb-appendix}
本附录补两条在正文被引用、但放正文会打断主线的恒等式证明。

\begin{derivation}[恒等式 (I5)：$(\mathbf{R}\mathbf{a})^\wedge=\mathbf{R}\mathbf{a}^\wedge\mathbf{R}^\top$]
对任意 $\mathbf{x}\in\mathbb{R}^3$，用外积在旋转下的协变性 $\mathbf{R}(\mathbf{a}\times\mathbf{b})=(\mathbf{R}\mathbf{a})\times(\mathbf{R}\mathbf{b})$（旋转保右手系与长度，故保外积），令 $\mathbf{b}=\mathbf{R}^\top\mathbf{x}$：
\[
  (\mathbf{R}\mathbf{a})^\wedge\mathbf{x}=(\mathbf{R}\mathbf{a})\times\mathbf{x}=(\mathbf{R}\mathbf{a})\times(\mathbf{R}\mathbf{R}^\top\mathbf{x})
  =\mathbf{R}\big(\mathbf{a}\times(\mathbf{R}^\top\mathbf{x})\big)=\mathbf{R}\mathbf{a}^\wedge\mathbf{R}^\top\mathbf{x}.
\]
对任意 $\mathbf{x}$ 成立，故 $(\mathbf{R}\mathbf{a})^\wedge=\mathbf{R}\mathbf{a}^\wedge\mathbf{R}^\top$（Barfoot 习题 7.3）。
\end{derivation}

\begin{derivation}[恒等式 (I6)：三重积 BAC--CAB]
逐分量展开 $\mathbf{a}\times(\mathbf{b}\times\mathbf{c})$ 的第一分量：
\[
  [\mathbf{a}\times(\mathbf{b}\times\mathbf{c})]_1=a_2(\mathbf{b}\times\mathbf{c})_3-a_3(\mathbf{b}\times\mathbf{c})_2
  =a_2(b_1c_2-b_2c_1)-a_3(b_3c_1-b_1c_3).
\]
整理并配出 $b_1(\mathbf{a}^\top\mathbf{c})-c_1(\mathbf{a}^\top\mathbf{b})$（加减 $a_1b_1c_1$）：
\[
  =b_1(a_1c_1+a_2c_2+a_3c_3)-c_1(a_1b_1+a_2b_2+a_3b_3)=b_1(\mathbf{a}^\top\mathbf{c})-c_1(\mathbf{a}^\top\mathbf{b}).
\]
其余两分量同理，故 $\mathbf{a}\times(\mathbf{b}\times\mathbf{c})=\mathbf{b}(\mathbf{a}^\top\mathbf{c})-\mathbf{c}(\mathbf{a}^\top\mathbf{b})$。等价矩阵形式 $\mathbf{a}^\wedge\mathbf{b}^\wedge=\mathbf{b}\mathbf{a}^\top-(\mathbf{a}^\top\mathbf{b})\mathbf{I}$。
\end{derivation}
